\documentclass[11pt,letterpaper]{article}
\usepackage[T1]{fontenc}
\usepackage[utf8]{inputenc}
\usepackage{lmodern}
\usepackage[margin=1in,headheight=15pt]{geometry}
\usepackage{amsmath,amssymb,amsthm,mathtools,mathrsfs}
\usepackage{microtype,booktabs,longtable,array,enumitem}
\usepackage[dvipsnames]{xcolor}
\usepackage{fancyhdr}
\usepackage[colorlinks=true,linkcolor=MidnightBlue,citecolor=MidnightBlue,urlcolor=MidnightBlue]{hyperref}
\usepackage[nameinlink,noabbrev]{cleveref}
\hypersetup{%
  pdftitle={Infinitesimal Analytic Geometry over the Surcomplex Numbers},
  pdfauthor={Merged research report},
  pdfsubject={Hahn analytic germ rings, Nullstellensatz, finite maps, residues}}
\pagestyle{fancy}\fancyhf{}
\fancyhead[L]{\small Infinitesimal surcomplex analytic geometry}
\fancyhead[R]{\small\thepage}
\renewcommand{\headrulewidth}{0.3pt}
\setlength{\parskip}{0.3em}
\setlength{\parindent}{1.3em}
\setlist{nosep,leftmargin=2em}
\setlength{\emergencystretch}{2em}
\allowdisplaybreaks[2]

\newtheorem{theorem}{Theorem}[section]
\newtheorem{lemma}[theorem]{Lemma}
\newtheorem{proposition}[theorem]{Proposition}
\newtheorem{corollary}[theorem]{Corollary}
\theoremstyle{definition}
\newtheorem{definition}[theorem]{Definition}
\newtheorem{example}[theorem]{Example}
\theoremstyle{remark}
\newtheorem{remark}[theorem]{Remark}
\newtheorem{convention}[theorem]{Convention}

\newcommand{\C}{\mathbb C}
\newcommand{\R}{\mathbb R}
\newcommand{\Q}{\mathbb Q}
\newcommand{\N}{\mathbb N}
\newcommand{\No}{\mathbf{No}}
\newcommand{\SC}{\mathbf{SC}}
\newcommand{\A}{\mathscr A}
\newcommand{\RR}{\mathscr R}
\newcommand{\FF}{\mathscr F}
\newcommand{\HH}{\mathcal H}
\newcommand{\m}{\mathfrak m}
\newcommand{\pp}{\mathfrak p}
\newcommand{\OO}{\mathcal O}
\newcommand{\Res}{\operatorname{Res}}
\newcommand{\res}{\operatorname{res}}
\newcommand{\supp}{\operatorname{supp}}
\newcommand{\lc}{\operatorname{lc}}
\newcommand{\id}{\operatorname{id}}
\newcommand{\Tr}{\operatorname{Tr}}
\newcommand{\ev}{\operatorname{ev}}
\newcommand{\Max}{\operatorname{Max}}
\newcommand{\Frac}{\operatorname{Frac}}
\newcommand{\disc}{\operatorname{disc}}
\newcommand{\NF}{\operatorname{NF}}
\newcommand{\red}{\operatorname{red}}
\newcommand{\vH}{v_{\!H}}
\newcommand{\vmin}{v_{\min}}
\newcommand{\HInt}{\mathop{\int}\nolimits^{\!H}}
\newcolumntype{L}[1]{>{\raggedright\arraybackslash}p{#1}}
\numberwithin{equation}{section}

\begin{document}

\hypersetup{pageanchor=false}
\begin{titlepage}
\centering
\vspace*{0.7cm}
{\Large\scshape A merged research report\par}
\vspace{0.7cm}
{\Huge\bfseries Infinitesimal Analytic Geometry\par}
\vspace{0.35cm}
{\LARGE over the Surcomplex Numbers\par}
\vspace{0.75cm}
{\Large Radius-Free Germs, a Nullstellensatz,\par and Finite Maps\par}
\vspace{0.65cm}
{\large September 21, 2026\par}
\vspace{0.5cm}
\begin{minipage}{0.93\textwidth}
\begin{abstract}
\noindent
We develop the several-variable analytic geometry of the infinitesimal
neighbourhood of an ordinary point in the surcomplex numbers
$\SC=\No[i]$, and we do so over \emph{two} coefficient rings whose
theorems look alike and are not interchangeable. Fix a nonzero divisible
set-sized ordered subgroup $\Gamma\subset\No$ and put $K=\C((t^\Gamma))$.
The \emph{radius-free} ring is
$\A_n=\C\{z_1,\ldots,z_n\}((t^\Gamma))$, whose Hahn coefficients are
individually convergent germs subject to no common radius; the
\emph{common-domain} ring $\RR_n$ consists of germs of Hahn families all of
whose coefficients are holomorphic on one ordinary neighbourhood. We prove
$\RR_n\subsetneq\A_n\subsetneq\C[[z]]((t^\Gamma))$, with explicit witnesses
in both directions, and we prove separately of each of $\A_n$ and $\RR_n$
that it is a Noetherian Jacobson domain of dimension $n$ whose maximal
ideals are exactly the infinitesimal evaluation ideals
$\m_a=(z_1-a_1,\ldots,z_n-a_n)$, $a\in\m_K^n$, with residue field $K$,
regular local rings, and completions $K[[X_1,\ldots,X_n]]$. Neither ring is
local: each records every infinitesimal point at once.

Over $\A_n$ we prove support-controlled Weierstrass division and
preparation, the strong Nullstellensatz $I(Z(I))=\sqrt I$ on the monad
$\m_K^n$ --- with an example showing that the nonzero-value-group
hypothesis is essential --- invariance of a finite zero scheme under
enlargement of the workspace, a full-surcomplex corollary accounting for
every infinitesimal zero in the proper class $\SC$, faithful flatness of the
inclusion into the formal-coefficient ring, and a sharp inverse-Jacobian
Hensel--Rouch\'e theorem with root matching. Over $\RR_n$ we prove
parameterized preparation and division on one fixed domain, Noether
normalization, finite surjective projections of irreducible monad zero sets
with generically reduced fibres of the expected degree, vanishing of all
higher homology of the Hahn Koszul complex, and a base-change theorem.
Both rings carry conservation of the finite length of an arbitrary isolated
complete intersection under positive-support perturbation, a perfect
residue pairing that survives collisions, and a trace--Jacobian identity;
the two residue functionals, one defined by a strongly summable series and
one by expansion on a fixed ordinary cycle, are shown to agree where both
are defined. Worked examples include a four-point system splitting across
different valuation ranks, a genuinely radius-free two-variable system with
infinite individual residues summing to an infinitesimal, and a
length-five collision with the exact collision equation
$3125u^6-108s^5=0$.
\end{abstract}
\end{minipage}
\vfill
\begin{minipage}{0.93\textwidth}\small
\textbf{Status.} Every assertion below is proved from explicitly identified
classical inputs. Multivariable Hahn--Weierstrass theory and
non-Archimedean Nullstellensatz results have substantial published
predecessors and are not claimed here as new general ideas. The proposed
contribution is the precise combination: the radius-free coefficient ring
with its structure package and workspace theorems, the common-domain
normal-form and residue package with explicit support certificates, and the
comparison between them. The literature search was made on 21 September
2026 and was targeted rather than exhaustive. No named published conjecture
is claimed solved, no exhaustive priority certification is asserted, and
none of the proofs has been formally verified. Symbolic checks validate the
displayed examples only; they are not machine verification of the general
proofs.
\end{minipage}
\end{titlepage}

\hypersetup{pageanchor=true}
\pagenumbering{roman}
\tableofcontents
\clearpage
\pagenumbering{arabic}

\section{The problem, the two rings, and the contribution}\label{analytic:sec:intro}

\subsection{From one moving zero to an analytic zero scheme}
\label{analytic:subsec:question}
Two supplied manuscripts on surcomplex analysis construct one-variable
preparation and moving-zero residue theory on infinitesimal thickenings of
ordinary complex domains, and both identify several-variable division and
finite mappings as their principal further targets
\cite{CompanionA,CompanionB}. The dependence of the present report on them is
motivational, not logical: none of their unproved targets is used here as a
theorem.

The several-variable question is not whether a simple root lifts. The
implicit function theorem already answers that. The substantial question is
whether a \emph{singular} system in several variables retains its entire
local algebra when its equations acquire arbitrary infinitesimal Hahn
perturbations. The ordinary system $(x^2,y^2)$ has a zero of length four; an
infinitesimal coupling can split it into four points lying at different
coordinate scales. Counting only the ordinary leading zero sees one point;
counting only distinct perturbed zeros misses collisions; and an argument
that treats all infinitesimals as powers of a single real-valued parameter
can miss entire valuation ranks. The appropriate invariant is the dimension
of the finite quotient algebra, together with its residue pairing and its
multiplication operators.

A second issue arises at the same moment, and it is the issue that makes
this report a report about two rings rather than one. Infinitely many
ordinary analytic coefficients may have radii of convergence tending to
zero. Requiring one common ordinary polydisk excludes them, even though
each coefficient has an unambiguous Taylor expansion at an infinitesimal
argument. Dropping the requirement admits them, but then no fixed ordinary
domain is available on which to run a Cartan-theoretic or contour argument.
Both choices are legitimate, both support a structure theory, and the
theorems they support are \emph{not} the same theorems, because they are
theorems about different rings.

\subsection{The four coefficient rings, and why they must be kept
apart}\label{analytic:subsec:four-rings}
Throughout, $\Gamma\subset\No$ is a nonzero divisible set-sized ordered
subgroup, $K=\C((t^\Gamma))$, $\m_K$ is its ideal of infinitesimals, and
$R_n=\C\{z_1,\ldots,z_n\}$ is the ring of convergent germs at the origin.
Four coefficient rings occur in this subject under closely similar
notation, and all four carry closely similar structural statements:
\begin{enumerate}
\item the \emph{fixed-polydisk} ring $\OO(D)((t^\Gamma))$, for one ordinary
polydisk $D$ held fixed once and for all;
\item the \emph{common-domain germ ring}
$\RR_n=\varinjlim_{U\ni0}\HH_\Gamma(U)$, where $\HH_\Gamma(U)$ consists of
Hahn families all of whose coefficient functions are holomorphic on the one
ordinary neighbourhood $U$, and the limit is over the restriction maps;
\item the \emph{radius-free} ring $\A_n=R_n((t^\Gamma))$, whose Hahn
coefficients are individually convergent germs with no common domain
required;
\item the \emph{formal-coefficient} ring $\FF_n=\C[[z_1,\ldots,z_n]]((t^\Gamma))$.
\end{enumerate}
There are natural maps
\begin{equation}\label{analytic:eq:ring-chain}
 \OO(D)((t^\Gamma))\longrightarrow\RR_n
 \subsetneq\A_n\subsetneq\FF_n ,
\end{equation}
and the two inclusions displayed as strict are proved strict in
\cref{analytic:ex:no-radius,analytic:rem:formal-strict}. Each of the four
rings satisfies, \emph{of itself}, a statement of the form ``the ring is
Noetherian; its maximal ideals are the infinitesimal evaluation ideals; its
local rings are regular with completion $K[[X]]$; division holds with
support in $S+E^*$.'' The statements are nearly word for word identical and
the notation nearly collides. This is a definitional divergence that reads
like a contradiction and is not one; it is also the single likeliest place
in this subject to produce a false attribution. We therefore adopt, and
keep, the following discipline.

\begin{convention}[Ring discipline]\label{analytic:conv:ring}
Every numbered statement below names the ring over which it is asserted. No
theorem proved of one of the four rings is transferred to another. Where a
statement holds of two rings, it is stated twice, with the two proofs given
and compared, because the proofs differ in an essential resource: over
$\A_n$ one may use ordinary germ-level division maps but has no fixed
ordinary domain; over $\RR_n$ one has a fixed ordinary domain and may use
Cartan's theorem~B and an ordinary residue cycle on it, but one cannot
reach a coefficient family with shrinking radii. We do not unify the
notation into a single symbol.
\end{convention}

The two results that relate the rings are the ones to invoke instead of a
transfer. The first is \cref{analytic:ex:no-radius}: the element
\[
 F(z)=\sum_{r\ge1}\frac{t^{r\eta}}{1-rz}\in\A_1
\]
has a pole at $1/r$ in its coefficient at $r\eta$, so it lies in no
$\HH_\Gamma(U)$ and is not an element of $\RR_1$ with those coefficients.
Consequently invoking a theorem stated only for coefficients holomorphic on
one fixed polydisk would leave a gap in the radius-free setting. The second
is \cref{analytic:thm:formal-comparison}: the inclusion $\A_n\hookrightarrow\FF_n$
is faithfully flat, with the same maximal ideals, residue fields,
completions and finite-colength quotients --- so the passage from
convergent to formal coefficients, unlike the passage from common-domain to
radius-free coefficients, changes nothing about a finite zero scheme.

\subsection{Overview of the main statements}\label{analytic:subsec:overview}
The following four theorems summarize the report. Each names its ring.

\begin{theorem}[Structure of the radius-free ring $\A_n$, overview]
\label{analytic:thm:overview-structure-A}
The ring $\A_n=\C\{z_1,\ldots,z_n\}((t^\Gamma))$ is a regular Noetherian
Jacobson domain of dimension $n$. There is a bijection
\[
 \m_K^n\longrightarrow\Max(\A_n),\qquad
 a\longmapsto\m_a=(z_1-a_1,\ldots,z_n-a_n),
\]
every residue field is $K$, for every ideal $I\subseteq\A_n$ one has
$I(Z_{\m_K^n}(I))=\sqrt I$, and the completion of $(\A_n)_{\m_a}$ is
$K[[z-a]]$. The ring $\A_n$ is not local.
\end{theorem}

\begin{theorem}[Structure and finite maps for the common-domain ring
$\RR_n$, overview]\label{analytic:thm:overview-structure-R}
The ring $\RR_n$ is a Noetherian Jacobson domain of dimension $n$ with the
same description of its maximal ideals, residue fields, local rings and
completions. In addition, every quotient $\RR_n/\pp$ by a prime of
dimension $d$ is finite over a coordinate copy of $\RR_d$; the induced
projection $V(\pp)\to\m_K^d$ of monad zero sets is surjective with finite
fibres, and over the complement of one nonzero $s\in\RR_d$ every fibre
consists of exactly $r=[\Frac(\RR_n/\pp):\Frac(\RR_d)]$ distinct reduced
points.
\end{theorem}

\begin{theorem}[Conservation, duality and the radius-free strengthening,
overview]\label{analytic:thm:overview-deform}
Let $f_1,\ldots,f_n\in R_n$ vanish at the origin with
$\mu=\dim_\C R_n/(f_1,\ldots,f_n)<\infty$, and let $b_1,\ldots,b_\mu$ be
monomials whose classes form a basis of $B_0=R_n/(f)$. Put $F_i=f_i+E_i$
with $\vH(E_i)>0$.
\begin{enumerate}
\item If the $E_i$ lie in $\RR_n$, then the classes of the $b_r$ are a
$K$-basis of $A_F=\RR_n/(F)$, all higher homology of the Hahn Koszul complex
of $F$ over $\RR_n$ vanishes, and all normal-form and division outputs on
one fixed ordinary domain have support in $\supp G+E^*$.
\item If the $E_i$ lie in the strictly larger ring $\A_n$, the same basis
statement and the same support certificate hold for
$B_F=\A_n/(F)$, by a different proof. In both cases $\dim_KB_F=\mu$, the
associated residue pairing is perfect, $\Res_F(HJ_F)=\Tr_K(M_H)$ and
$\Res_F(J_F)=\mu$, the zero set $Z_F\subset\m_K^n$ is finite and nonempty
with $\sum_{a\in Z_F}e_a=\mu$, and enlarging the workspace $\Gamma$ changes
neither the zeros nor their multiplicities.
\end{enumerate}
\end{theorem}

\begin{theorem}[Quantitative stability over $\A_n$, overview]
\label{analytic:thm:overview-stability}
Let $F\in\A_n^n$ have $\vH(F_i)\ge0$ and let $a\in\m_K^n$ be a simple zero.
Put $\kappa_a=\max\bigl(0,-\min_{i,j}v((DF(a)^{-1})_{ij})\bigr)$. If $G=F+D$
and $\delta=\min_i\vH(D_i)>2\kappa_a$, there is exactly one zero $b$ of $G$
with $\vmin(b-a)>\kappa_a$; it is simple and satisfies
$\vmin(b-a)\ge\delta-\kappa_a$. The strict factor two cannot be weakened.
\end{theorem}

\Cref{analytic:thm:overview-structure-A} is proved in
\cref{analytic:sec:division,analytic:sec:structure-A};
\cref{analytic:thm:overview-structure-R} in
\cref{analytic:sec:structure-R,analytic:sec:finite};
\cref{analytic:thm:overview-deform} in
\cref{analytic:sec:deform,analytic:sec:trace,analytic:sec:extension}; and
\cref{analytic:thm:overview-stability} in \cref{analytic:sec:stability}.
The maximal-ideal assertions are stronger than the mere ability to
substitute infinitesimals: they say that substitution accounts for
\emph{all} closed points of these algebras.

\subsection{Predecessors and the status of novelty}
\label{analytic:subsec:novelty}
A literature check changes the appropriate framing of these results, and we
adopt the more careful of the two framings available to us.

Cluckers--Lipshitz--Robinson prove multivariable preparation and division
for nonstandard analytic structures \cite[Theorem 2.5 and Remark
2.6]{CLR}. Cluckers--Lipshitz explicitly treat analytic coefficient
families with a \emph{common radius} and arbitrary well-ordered Hahn
supports \cite[\S3.1, Remark 3.1.4, Example 3.3(6)]{CL}, and their
separated analytic theory contains normalization and Nullstellensatz
theorems \cite[\S5.2]{CL}. These are close and substantial predecessors.
Passing to complex coefficients, phrasing their consequences in the
terminology of the supplied manuscripts, or giving another Weierstrass
recursion is not by itself a novelty claim. Neither a Hahn Weierstrass
theorem nor a non-Archimedean Nullstellensatz is claimed here as a new
general idea. The coefficient ring and its support restrictions must be
compared, not just the theorem titles: our radius-free ring $\A_n$ retains
convergence of each complex coefficient germ but imposes no common radius,
and \cref{analytic:ex:no-radius} proves that it strictly contains the
common-domain ring $\RR_n$. We do not claim that $\A_n$ lies outside every
broader analytic-structure framework.

Conservation of intersection number and residue duality are classical in
complex analytic geometry \cite{GunningRossi,GH,SS,GrauertRemmert}, and the
perturbation mechanism used over $\RR_n$ is a special case of homological
perturbation theory \cite{Crainic}. Residue expansions, normal forms and
trace identities have classical antecedents, notably
Cattani--Dickenstein--Sturmfels \cite{CDS}, Scheja--Storch \cite{SS} and
Kunz--Waldi \cite{KW}. Surreal analysis likewise has substantial
antecedents: Alling develops analytic constructions over surreal number
fields \cite{Alling}, Berarducci--Mantova study transseries as surreal
function germs \cite{BM}, and Costin--Ehrlich develop surreal extension and
integration theory \cite{CE}. Our differentiation is differentiation in the
variables $z_i$ with Hahn coefficients held fixed; it is not a derivation on
the field of surreal numbers.

The new-work claim is deliberately narrow. We give an explicit
fixed-neighbourhood realization of the perturbation mechanism for arbitrary
Hahn supports, retain a specified basis throughout the deformation, obtain
uniform support certificates for normal forms and matrices, prove
compatibility with all divisible workspace enlargements, identify the
associated moving-root residue functional including its collision-stable
pairing, and --- separately, over the strictly larger radius-free ring ---
prove the structure package, the strong Nullstellensatz, the workspace and
full-surcomplex theorems, the analytic--formal comparison, and the
quantitative root correspondence. In the sources inspected we did not
locate these exact combined statements for either ring. We offer them as
proposed new results in a precisely specified setting; they may also be
derivable from a broader existing framework. This is a proposed original
synthesis and formulation, not a certified first theorem in multivariable
non-Archimedean analysis. The search was conducted on 21 September 2026 and
was not a systematic review of all publications. The proofs, rather than a
priority assertion, are the substantive deliverable.

\subsection{Notation fixed once}\label{analytic:subsec:conventions}
The sources of this report disagree on several symbols and on one name. We
fix the following and use nothing else.

\begin{convention}[Notation]\label{analytic:conv:notation}
\begin{enumerate}
\item \emph{Reduction.} The map $K^\circ\to\C$ extracting the coefficient of
$t^0$ is written $\red$, and the associated topology on $\m_K^n$ obtained by
pulling back the ordinary topology of $\C^n$ along standard parts is called
the \emph{reduction topology}. It is called the residue topology and the
shadow topology elsewhere; those are synonyms for the same object and are
not used again.
\item \emph{Evaluation.} The value of $F$ at an infinitesimal point $a$ is
written $F(a)$. The notation $F^\sharp(a)$ occurs in the sources for the
same operation and is not used again.
\item \emph{Valuations.} $v$ is the field valuation on $K$, with
$v(0)=+\infty$; $\vH$ is the Hahn valuation $\min\supp$ on a ring of Hahn
series with coefficients in a coefficient ring; $\vmin(a)=\min_iv(a_i)$ for
a tuple.
\item \emph{Support monoid.} For well-ordered $E\subset\Gamma_{>0}$ we write
$E^*=\{0\}\cup\{e_1+\cdots+e_r:r\ge1,\ e_j\in E\}$. The symbol $T^*$ is used
for the same object in one source and is not used again.
\item \emph{Deformation data.} $\mu$ is the ordinary colength;
$b_1,\ldots,b_\mu$ the fixed monomial basis; $\NF_F$ the normal-form map
with coordinates $\NF_F(G)_r$; $Q_j$ the division certificates; $M_H$ the
multiplication matrix of $H$; $\mathcal G_F$ the residue Gram matrix; $e_a$
the local algebraic multiplicity at $a$ (written $\mu_a(F)$ in one source);
$J_F$ the Jacobian determinant; $\kappa_a$ the inverse-Jacobian valuation
loss.
\item \emph{Residues.} $\res_g$ is the ordinary local Grothendieck residue.
$\Res_F$ is the \emph{series} functional \eqref{analytic:eq:hahn-residue},
defined over $\A_n$; $\Lambda_F$ is the \emph{contour-expansion} functional
\eqref{analytic:eq:hahn-contour-residue}, defined over $\RR_n$. They are
distinct definitions and are shown to agree where both apply
(\cref{analytic:prop:residues-agree}).
\item $\N$ contains zero. $\SC=\No[i]$, and $t^\gamma=\omega^{-\gamma}$.
\end{enumerate}
\end{convention}

\section{Hahn workspaces and the summability mechanism}
\label{analytic:sec:hahn}

This section states the foundations that recur in every later proof, in the
one place where they belong. Everything in it is common to both coefficient
rings; it is the last material in this report of which that is true.

\subsection{Set-sized workspaces inside the surcomplex class}
\label{analytic:subsec:workspace}
A surcomplex number is an element of $\SC=\No[i]$. We use the normal-form
monomials $t^\gamma=\omega^{-\gamma}$, so that a positive exponent means an
infinitesimal. The Conway normal-form theorem identifies $\C((t^\Gamma))$
with a subfield of $\SC$: its elements have well-ordered exponent support
and complex coefficients. The Hahn-field theorem says that
$K=K_\Gamma=\C((t^\Gamma))$ is algebraically closed when $\Gamma$ is
divisible and the coefficient field is algebraically closed of
characteristic zero; see \cite[Corollary 4]{Poonen}. Write
\[
 R_\Gamma=\R((t^\Gamma)),\qquad K=R_\Gamma[i],\qquad
 K^\circ=\{a:v(a)\ge0\},\qquad
 \m_K=\{a\in K:v(a)>0\}\cup\{0\},
\]
with $R_\Gamma$ real closed, $\red:K^\circ\to\C$ the reduction of
\cref{analytic:conv:notation}, and $\vmin(a)=\min_iv(a_i)$ for tuples. Under
the surreal embedding, $\m_K$ consists exactly of the elements of $K$ that
are infinitesimal relative to the positive real numbers.

\begin{lemma}[Enlargement rule]\label{analytic:lem:enlargement}
Every finite tuple of surcomplex numbers lies in some set-sized
algebraically closed workspace $K_\Delta$ of the above form: take the
additive subgroup generated by the union of the normal-form supports of the
entries, enlarge by a positive element if necessary, and pass to the
divisible hull.
\end{lemma}

This observation is what later turns a workspace theorem into a statement
about all surcomplex zeros. Its limit must be stated with equal force: it
never authorizes a proper-class sum or a proper-class polynomial ring. No
proper-class construction is used anywhere below. Every theorem about a ring
concerns a fixed set-sized $\Gamma$, and the passage to the full class is
made only by the explicit base-change arguments of
\cref{analytic:sec:extension}, for one finite cluster at a time.

\subsection{Strong summability and the support calculus}
\label{analytic:subsec:summability}
\begin{definition}[Strong summability and the monoid $E^*$]
\label{analytic:def:strong-sum}
A family of Hahn series is \emph{strongly summable} if the union of its
supports is well ordered and only finitely many members contribute to each
exponent; its sum is then defined coefficientwise. For a well-ordered
$E\subset\Gamma_{>0}$,
\[
 E^*=\{0\}\cup\{e_1+\cdots+e_r:r\ge1,\ e_j\in E\}.
\]
\end{definition}

This is the only notion of infinite sum used in this report. It replaces
any topological convergence, and no continuity or norm hypothesis is ever
invoked in its place.

\begin{lemma}[Hahn--Neumann support lemma]\label{analytic:lem:neumann}
Let $S_1,S_2$ be well-ordered subsets of an ordered abelian group. Then
$S_1+S_2$ is well ordered and every exponent has only finitely many
representations $s_1+s_2$ with $s_i\in S_i$. If $E\subset\Gamma_{>0}$ is
well ordered then $E^*$ is well ordered, and a fixed exponent has only
finitely many representations as a finite word in $E$, including finitely
many possible word lengths. Consequently, for well-ordered $S$, the data
indexed by a member of $S$ together with a word in $E$ are locally finite
over $S+E^*$.
\end{lemma}

The empty word represents zero. Positivity of $E$ is essential: zero
increments would permit infinitely many contributions at one exponent. These
facts hold verbatim for supports of vector-valued series. The ring
construction and the geometric inversion are reviewed in \cite[\S3]{Poonen};
the ordered-series support theorem originates in \cite{Neumann}.

\begin{definition}[Hahn families of vectors]\label{analytic:def:vector-hahn}
For a complex vector space $V$, let $V((t^\Gamma))$ denote the well-ordered
Hahn families with coefficients in $V$. A complex-linear map $L:V\to W$ acts
coefficientwise and never increases support.
\end{definition}

The last sentence carries more weight than its length suggests. It makes no
continuity assumption on $L$, so ordinary holomorphic division maps may be
applied to infinitely many coefficient functions even when those functions
have arbitrarily large ordinary norms, and --- over $\A_n$ --- even when
they share no analytic neighbourhood at all.

\begin{lemma}[Positive-support operator inversion]
\label{analytic:lem:operator}
Let $V$ be a complex vector space, let $E\subset\Gamma_{>0}$ be well
ordered, and let
\[
 T=\sum_{e\in E}t^eT_e
\]
act on $V((t^\Gamma))$, each $T_e:V\to V$ complex-linear. Then
\begin{equation}\label{analytic:eq:operator-inverse}
 (\id+T)^{-1}G=\sum_{r\ge0}(-T)^rG
\end{equation}
is strongly summable for every $G$, and defines a $K$-linear operator. If
$S=\supp G$, the output support lies in $S+E^*$. The same statement holds
for finite direct sums of vector spaces and for matrices of such operators.
\end{lemma}
\begin{proof}
A contribution to the coefficient at $\nu$ in $T^rG$ is indexed by
$s+e_1+\cdots+e_r=\nu$ with $s\in S$ and $e_j\in E$. By
\cref{analytic:lem:neumann} there are only finitely many such words,
including finitely many possible lengths, so each output coefficient is a
finite sum of compositions of the $T_e$ applied to finitely many input
coefficients. The union of the possible exponents lies in the well-ordered
set $S+E^*$. Multiplication by $\id+T$ telescopes coefficientwise to $G$ on
either side. Commuting an additional scalar Hahn series through these
operators is justified by the same finite-contribution property, which
gives $K$-linearity. A finite matrix index set changes nothing.
\end{proof}

\subsection{The two negative companions}\label{analytic:subsec:negatives}
Both facts in this subsection are load-bearing. Omitting either invites a
proof that looks like an iteration argument and is not one.

\begin{remark}[Strong summability is not a valuation limit]
\label{analytic:rem:rank}
Take $\Gamma=\Q+\Q\omega$ with its inherited order; for every natural
number $N$ one has $N<\omega$. The partial sums of $\sum_{r\ge0}t^r$
therefore need not converge in the valuation topology: their tails never
have valuation greater than $\omega$. Nevertheless the family is strongly
summable and its sum is $(1-t)^{-1}$. The same phenomenon appears in
$\Q\oplus\Q$ ordered lexicographically: with $e=(0,1)$ and $b=(1,0)$, every
$ne$ is below $b$, yet $\sum_{n\ge0}t^{ne}=(1-t^e)^{-1}$ while the errors
never have valuation above $b$. Hence no argument of the form ``the
valuation improves by a fixed positive amount at each iteration'' justifies
any construction in this report. \Cref{analytic:lem:operator} does justify
them, and replacing it by an assertion that errors tend to zero would
discard legitimate cases of every theorem below.
\end{remark}

\begin{remark}[Finitely many words at one exponent is not finitely many
exponents below it]\label{analytic:rem:initial-segments}
A well-ordered positive support may have infinitely many exponents below a
prescribed bound. Our proofs never require those initial segments to be
finite; they require finitely many \emph{words} contributing to one exact
exponent. For $E=\{2-1/m:m\ge1\}$ there are infinitely many members below
$2$, yet the coefficient of a given Hahn product is still a finite
convolution. The distinction is precisely what makes the
finite-parameter-replacement arguments of
\cref{analytic:lem:finite-dependency} and
\cref{analytic:thm:trace} legitimate: at a prescribed exponent, only
finitely many coefficient germs are involved, and that is all a comparison
with an ordinary finite-parameter deformation needs. In $\Q+\Q\omega$ the
assertion ``only finitely many exponents below'' is false below $\omega$.
\end{remark}

\begin{remark}[Well-ordering cannot be weakened to positivity]
\label{analytic:rem:positivity}
The set $\{1/r:r\ge1\}$ is positive but not well ordered, so the
corresponding putative Hahn series is not an element of any of the four
rings of \cref{analytic:subsec:four-rings}. Positivity alone would break
\cref{analytic:lem:neumann} and with it every support certificate below.
\end{remark}

\section{The coefficient rings and their values}\label{analytic:sec:germs}

\subsection{The radius-free ring $\A_n$}\label{analytic:subsec:ringA}
\begin{definition}[The radius-free Hahn-germ ring]\label{analytic:def:ringA}
Let $R_n=\C\{z_1,\ldots,z_n\}$ and define
\begin{equation}\label{analytic:eq:ringA}
 \A_n=R_n((t^\Gamma))
 =\left\{\sum_{\gamma\in S}f_\gamma(z)t^\gamma:
 S\subset\Gamma\ \text{well ordered},\ f_\gamma\in R_n\right\}.
\end{equation}
For $F\ne0$ put $\vH(F)=\min\supp F$ and let $\lc(F)$ be its leading
coefficient germ. A hypothesis $\vH(F)>0$ permits $F=0$.
\end{definition}

The symbol $\A_n$ does not mean an unspecified ring of convergent series
over $K$: \cref{analytic:def:ringA} is the entire convention. In particular
no convergence condition on the numerical sizes of the complex Taylor
coefficients is imposed uniformly in $\gamma$, and the radii of the
$f_\gamma$ may tend to zero. The Hahn construction makes $\A_n$ a
$K$-algebra and an integral domain, since the leading coefficient of a
product is the product of the leading coefficient germs in the domain
$R_n$; and $\A_0=K$.

\begin{proposition}[Units of $\A_n$]\label{analytic:prop:units}
For $F=t^\lambda(f_0+E)\in\A_n$ with $\vH(E)>0$ and $f_0\ne0$, the
following are equivalent: $F$ is a unit of $\A_n$; $f_0$ is a unit of $R_n$;
$f_0(0)\ne0$.
\end{proposition}
\begin{proof}
If $f_0$ is a unit then
$F^{-1}=t^{-\lambda}f_0^{-1}\sum_{r\ge0}(-f_0^{-1}E)^r$ is strongly
summable by \cref{analytic:lem:operator}. Conversely if $FG=1$ then the
product of the leading coefficient germs is $1$, so $f_0$ is a unit. The
last equivalence is ordinary local analytic algebra.
\end{proof}

\begin{proposition}[Evaluation and translation over $\A_n$]
\label{analytic:prop:translation}
For $a\in\m_K^n$ and $F=\sum f_\gamma t^\gamma\in\A_n$, the expression
\begin{equation}\label{analytic:eq:evaluation}
 F(a)=\sum_{\gamma,\alpha}
 \frac{\partial^\alpha f_\gamma(0)}{\alpha!}\,t^\gamma a^\alpha
\end{equation}
is strongly summable and defines a $K$-algebra homomorphism
$\ev_a:\A_n\to K$. The formula
\begin{equation}\label{analytic:eq:translation}
 \tau_aF(z)=\sum_{\gamma,\alpha}
 \frac{\partial^\alpha f_\gamma(z)}{\alpha!}\,t^\gamma a^\alpha
\end{equation}
defines an automorphism of $\A_n$ with $\tau_a\tau_b=\tau_{a+b}$,
$\tau_a^{-1}=\tau_{-a}$ and $(\tau_aF)(b)=F(a+b)$.
\end{proposition}
\begin{proof}
Let $S=\supp F$ and let $E_a$ be the union of the supports of the
coordinates of $a$, a well-ordered positive set. Every expanded term in
either formula has support in $S+E_a^*$. At a fixed exponent,
\cref{analytic:lem:neumann} permits only finitely many source exponents,
coordinate choices, and products of coefficients of $a$; hence each
coefficient in \eqref{analytic:eq:translation} is a finite sum of
derivatives of ordinary analytic germs and is again an analytic germ. The
same argument gives strong summability of
\eqref{analytic:eq:evaluation}. The product rule, the multinomial identity
and the formal Taylor translation identity then give all the asserted
algebraic identities, the rearrangements being finite at each exponent.
Note what this proof does \emph{not} do: it never translates an ordinary
germ by a nonzero complex number outside its domain. It translates by an
infinitesimal, using derivatives at the original centre.
\end{proof}

\begin{proposition}[Evaluation kernels and finite jets over $\A_n$]
\label{analytic:prop:kernels}
For every $a\in\m_K^n$,
\begin{equation}\label{analytic:eq:kernel}
 \ker\ev_a=\m_a=(z_1-a_1,\ldots,z_n-a_n),
\end{equation}
and for $N\ge1$ the Taylor map induces
\begin{equation}\label{analytic:eq:jets}
 \A_n/\m_a^N\simeq K[X_1,\ldots,X_n]/(X_1,\ldots,X_n)^N.
\end{equation}
\end{proposition}
\begin{proof}
At $a=0$ apply ordinary Taylor division to each coefficient germ: if its
constant term vanishes, write it as $\sum_iz_ig_i$ using fixed
complex-linear coordinate-division maps, and extend those maps
coefficientwise. For general $N$, subtract the Taylor polynomial of degree
below $N$ and divide the remainder by the finitely many monomials of degree
$N$; polynomial representatives give surjectivity. Translation by
\cref{analytic:prop:translation} reduces arbitrary $a$ to zero.
\end{proof}

\begin{proposition}[Faithful interpretation as functions on the monad]
\label{analytic:prop:faithful}
If $F\in\A_n$ vanishes at every point of $\m_K^n$ then $F=0$.
\end{proposition}
\begin{proof}
The case $n=0$ is immediate. Normalize a nonzero $F$ as $f_0+E$ and let $m$
be the ordinary vanishing order of $f_0$. Choose $c\in\C^n$ at which the
degree-$m$ homogeneous part $H_m$ of $f_0$ is nonzero. If $E\ne0$ set
$e=\vH(E)>0$ and choose $\eta>0$ in $\Gamma$ with $m\eta<e$, which
divisibility permits (for instance $\eta=e/(m+1)$); if $E=0$ take any
$\eta>0$. At $a=ct^\eta$ one has
$f_0(a)=H_m(c)t^{m\eta}+(\text{larger valuation})$ while $v(E(a))\ge e$, so
$F(a)\ne0$. The case $m=0$ is included.
\end{proof}

\begin{remark}[Formal power series are different]
\label{analytic:rem:formal-distinction}
The Taylor embedding $\A_n\hookrightarrow K[[z]]$ is not an identification
of rings of functions. The coefficientwise formal inverse
\[
 (z-t^\eta)^{-1}=-\sum_{r\ge0}t^{-(r+1)\eta}z^r
\]
lies in $K[[z]]$ but not in $\A_1$, since its combined Hahn support is not
well ordered. Accordingly $z-t^\eta$ is a nonunit of $\A_1$, and it actually
vanishes at an infinitesimal point. Completing at a point and working on the
entire monad are different operations.
\end{remark}

\subsection{The common-domain ring $\RR_n$}\label{analytic:subsec:ringR}
\begin{definition}[Fixed-domain Hahn families and the common-domain germ
ring]\label{analytic:def:ringR}
For an ordinary connected open $U\subset\C^n$ set
\[
 \HH_\Gamma(U)=\Bigl\{\sum_{\gamma\in S}f_\gamma t^\gamma:
 f_\gamma\in\OO(U),\ S\subset\Gamma\ \text{well ordered}\Bigr\},
\]
and define the common-domain germ ring at the ordinary origin by
\begin{equation}\label{analytic:eq:ringR}
 \RR_n=\RR_n(\Gamma)=\varinjlim_{U\ni0}\HH_\Gamma(U),
\end{equation}
the limit being over connected ordinary neighbourhoods of $0$ with the
restriction maps. Set $\RR_0=K$. For $c\in U$ and $\varepsilon\in\m_K^n$ the
evaluation of $F=\sum f_\gamma t^\gamma$ is
\begin{equation}\label{analytic:eq:evaluationR}
 F(c+\varepsilon)=\sum_{\gamma\in\supp F,\ \alpha\in\N^n}
 t^\gamma\frac{\partial^\alpha f_\gamma(c)}{\alpha!}\varepsilon^\alpha .
\end{equation}
\end{definition}

One may restrict the limit in \eqref{analytic:eq:ringR} to ordinary
polydisks. Restriction is injective by the ordinary identity theorem, and
multiplication of leading coefficient germs shows $\RR_n$ is a domain. Every
element of $\RR_n$ evaluates on all of $\m_K^n$; its elements are functions
on the \emph{whole} monad, not germs at one point of a fine topology.

The quantifier order in \eqref{analytic:eq:ringR} is the entire point. All
coefficient functions of one representative have \emph{one} common domain.
The larger expression $\OO_{\C^n,0}((t^\Gamma))=\A_n$ permits independently
shrinking coefficient domains and is a different ring.

\begin{proposition}[Evaluation, support and translation over $\RR_n$]
\label{analytic:prop:evaluationR}
Formula \eqref{analytic:eq:evaluationR} is strongly summable; evaluation
respects sums, products and coefficientwise partial derivatives, and is
injective on $\HH_\Gamma(U)$. Composition with an infinitesimal translation,
or with a positive-support perturbation of an ordinary holomorphic map, is
well defined whenever the ordinary domains match. For $a\in\m_K^n$,
translation $\tau_aF(z)=F(z+a)$ is an automorphism of $\RR_n$ with inverse
$\tau_{-a}$, and
\begin{equation}\label{analytic:eq:coordinatequotient}
 \RR_n/(z_1-a_1,\ldots,z_k-a_k)\simeq\RR_{n-k}\qquad(0\le k\le n).
\end{equation}
In particular $\ev_a:\RR_n\to K$ is surjective with kernel
$\m_a=(z_1-a_1,\ldots,z_n-a_n)$.
\end{proposition}
\begin{proof}
Put $P=\bigcup_j\supp\varepsilon_j$. All exponents in
\eqref{analytic:eq:evaluationR} lie in $\supp F+P^*$, and the multiindices
of each fixed total degree give finitely many choices, so
\cref{analytic:lem:neumann} gives local finiteness and the ordinary Taylor
product and composition identities may be rearranged coefficientwise.
Differentiating coefficient functions introduces no new exponents. At
ordinary points $F(c)=\sum f_\gamma(c)t^\gamma$; if all these vanish, then
uniqueness of Hahn coefficients makes every $f_\gamma$ identically zero,
which is injectivity. For translation, the Taylor formula has support
bounded by $\supp F$ plus the monoid generated by the coordinates of $a$,
and each derivative coefficient remains holomorphic on the same ordinary
domain; the binomial identity gives $\tau_a\tau_{-a}=\id$. For
\eqref{analytic:eq:coordinatequotient} at $a=0$, the ordinary identity
\[
 h(z_1,z')=h(0,z')+z_1\,\frac{h(z_1,z')-h(0,z')}{z_1}
\]
has a holomorphic quotient on any polydisk; applying it coefficientwise
eliminates one coordinate without changing support, and iterating
eliminates $k$. Translation handles general $a$.
\end{proof}

\begin{lemma}[Infinitesimal evaluation detects nonzero elements of $\RR_n$]
\label{analytic:lem:density}
For every nonzero $F\in\RR_n$ there is $a\in\m_K^n$ with $F(a)\ne0$.
\end{lemma}
\begin{proof}
Identical in mechanism to \cref{analytic:prop:faithful}, and we record it
separately because it is a statement about a different ring. Normalize
$F=t^\lambda(f_0+E)$ with $f_0\ne0$ and $\vH(E)>0$, let $m$ be the least
total degree of a nonzero homogeneous term of $f_0$, and choose $c\in\C^n$
at which that homogeneous polynomial is nonzero. If $E\ne0$ put
$e=\min\supp E>0$ and $\delta=e/(m+1)\in\Gamma_{>0}$, so $m\delta<e$;
otherwise take any $\delta>0$. At $a=t^\delta c$ the leading term of
$f_0(a)$ has valuation $m\delta$, while every value contributed by $E$ has
valuation at least $e$ and cannot cancel it.
\end{proof}

The nontriviality of $\Gamma$ is necessary here: when $\Gamma=\{0\}$ the
monad has only its origin and cannot detect ordinary analytic germs. This
is the same hypothesis whose necessity is made precise in
\cref{analytic:rem:trivial-group}.

\subsection{The two inclusions are strict}
\label{analytic:subsec:strictness}
\begin{example}[No shared ordinary neighbourhood: $\RR_n\subsetneq\A_n$]
\label{analytic:ex:no-radius}
Fix $\eta\in\Gamma_{>0}$ and consider
\begin{equation}\label{analytic:eq:no-radius}
 F(z)=\sum_{r\ge1}\frac{t^{r\eta}}{1-rz}.
\end{equation}
This is an element of $\A_1$ and it is evaluable at every infinitesimal
argument. Yet it cannot be represented by a Hahn family of holomorphic
functions on any common ordinary disk centred at zero, so it is not an
element of $\RR_1$ with those coefficients: its coefficient at $r\eta$ has a
pole at $1/r$, and an ordinary disk of positive radius contains $1/r$ for
all sufficiently large $r$. Equality of germwise Hahn expressions forces
equality of the coefficient germs, and so does equality as functions on the
whole monad, by \cref{analytic:prop:faithful}. Hence no representative with
a common domain exists.
\end{example}

Taking Hahn series and passing to germs therefore do not commute in the
naive common-radius sense. This is exactly why invoking a theorem stated
only for coefficients holomorphic on one fixed polydisk would leave a gap in
the radius-free setting, and it is the reason
\cref{analytic:conv:ring} is enforced below without exception.

\begin{remark}[$\A_n\subsetneq\FF_n$]\label{analytic:rem:formal-strict}
The constant-Hahn-support expression $\sum_{r\ge0}r!\,z^r$ lies in
$\FF_1=\C[[z]]((t^\Gamma))$ but not in $\A_1$, since $\sum r!z^r$ is not a
convergent germ. Both inclusions in
\eqref{analytic:eq:ring-chain} are therefore strict, and the comparison
theorem \cref{analytic:thm:formal-comparison} has content beyond a change
of notation.
\end{remark}

\subsection{Taylor calculus on the monad}\label{analytic:subsec:taylor}
Differentiation acts coefficientwise on all four rings. Translation gives a
Taylor expansion at every infinitesimal centre with a useful valuation
bound: if $\vH(F)\ge\lambda$, $a,h\in\m_K^n$ and $\vmin(h)=\rho>0$, then
\begin{equation}\label{analytic:eq:taylor-bound}
 v\Bigl(F(a+h)-\sum_{|\alpha|<N}
 \frac{\partial^\alpha F(a)}{\alpha!}h^\alpha\Bigr)\ \ge\ \lambda+N\rho,
\end{equation}
because translation does not decrease the lower support bound and every
remaining monomial contains at least $N$ factors from $h$. This gives the
expected differential and chain rules on the monad, and the estimates
persist after enlarging the workspace. Strong summability remains the
definition of the series: these estimates do not turn every sequence of
partial sums into a convergent sequence.

\section{Weierstrass division and preparation, over each ring}
\label{analytic:sec:division}

Both rings admit division with the support certificate $S+E^*$, and the two
proofs use different ordinary inputs. Over $\A_n$ the input is the ordinary
\emph{germ-level} division pair, which is available for each coefficient
separately and produces no common domain. Over $\RR_n$ the input is a pair
of \emph{contour} operators attached once to one product neighbourhood,
which never changes the domain. Neither proof can be run in the other
setting, and the two theorems are therefore stated separately.

\subsection{Division over the radius-free ring $\A_n$}
\label{analytic:subsec:divA}
Write $z=(x,w)$ with $x=(z_1,\ldots,z_{n-1})$ and $w=z_n$. An ordinary germ
$f\in R_n$ is \emph{$w$-regular of order $m$} if
$f(0,w)=cw^m+O(w^{m+1})$ with $c\ne0$. The classical Weierstrass theorem
gives $f=u_0P_0$ with $u_0$ a unit and
$P_0=w^m+\sum_{j<m}p_{0j}(x)w^j$, $p_{0j}(0)=0$; classical division gives
\emph{unique} $R_{n-1}$-linear maps
\[
 D:R_n\to R_n,\qquad \Pi:R_n\to R_{n-1}[w]_{<m},\qquad
 g=f\,Dg+\Pi g,
\]
so that $D(fg)=g$ and $D(r)=0$ when $\deg_wr<m$. Standard treatments include
\cite{GrauertRemmert}. Hahn analogues occur in \cite{CL,CLR}; the proof
below verifies directly the coefficient and support conditions we need.

\begin{theorem}[Hahn-germ division over $\A_n$]
\label{analytic:thm:divisionA}
Let $F=f+E\in\A_n$, where $f\in R_n$ is $w$-regular of order $m\ge1$ and
$\vH(E)>0$. Every $G\in\A_n$ has a unique decomposition
\begin{equation}\label{analytic:eq:divisionA}
 G=FQ+R,\qquad Q\in\A_n,\quad R\in\A_{n-1}[w]_{<m}.
\end{equation}
Writing $E_+=\supp E$ and $S=\supp G$, both $Q$ and $R$ have support
contained in $S+E_+^*$. Explicitly
\begin{equation}\label{analytic:eq:divisionA-formula}
 Q=(\id+DM_E)^{-1}DG,\qquad R=\Pi(G-EQ),
\end{equation}
where $M_E$ is multiplication by $E$ and $D,\Pi$ act coefficientwise.
\end{theorem}
\begin{proof}
Apply $D$ to the desired equation. Since $D(fQ)=Q$ and $D(R)=0$ it becomes
$(\id+DM_E)Q=DG$. The operator $DM_E$ is a sum of positive-support
complex-linear operators on $R_n$, so \cref{analytic:lem:operator} produces
the first formula in \eqref{analytic:eq:divisionA-formula} with the stated
support bound. That identity says $D(G-EQ)=Q$; applying ordinary division to
$G-EQ$ gives $G-EQ=fQ+\Pi(G-EQ)$, which is \eqref{analytic:eq:divisionA}.
The remainder has the required degree because $\Pi$ does so coefficientwise.
For uniqueness, subtract two decompositions and apply $D$: the difference of
the quotients lies in the kernel of the invertible operator $\id+DM_E$,
hence is zero, and then so is the difference of remainders. All produced
coefficients are ordinary analytic germs, since a fixed Hahn coefficient
involves only finitely many applications of ordinary division.
\end{proof}

\begin{theorem}[Preparation and finite projection over $\A_n$]
\label{analytic:thm:prepA}
Under the hypotheses of \cref{analytic:thm:divisionA} there is a unique
factorization
\begin{equation}\label{analytic:eq:prepA}
 F=UP,\qquad P=w^m+\sum_{j<m}p_j(x)w^j,
\end{equation}
with $U\in\A_n$ a unit, each $p_j\in\A_{n-1}$ of nonnegative support,
$\lc(P)=P_0$ and $\lc(U)=u_0$. The supports of $P$ and $U$ lie in $E_+^*$.
In particular
\begin{equation}\label{analytic:eq:free-hypersurfaceA}
 \A_n/(F)\simeq\A_{n-1}[w]/(P)
\end{equation}
is free of rank $m$ over $\A_{n-1}$, with basis $1,w,\ldots,w^{m-1}$.
\end{theorem}
\begin{proof}
Divide $w^m$ by $F$: $w^m=FQ+R$. The leading coefficient of $Q$ is
$D(w^m)=u_0^{-1}$, because $w^m=P_0+(w^m-P_0)$ is ordinary division by
$P_0$; hence $Q$ is a unit by \cref{analytic:prop:units}. Put $P=w^m-R$ and
$U=Q^{-1}$; their leading coefficients are $P_0$ and $u_0$, the support
assertion for $Q,R$ comes from division, and geometric inversion gives it
for $U$. If $F=\widetilde U\widetilde P$ is another such factorization then
$w^m=F\widetilde U^{-1}+(w^m-\widetilde P)$ is another division of $w^m$
with remainder of degree below $m$, so uniqueness in
\cref{analytic:thm:divisionA} forces $\widetilde P=P$ and
$\widetilde U=U$. Division gives both surjectivity and uniqueness of
polynomial representatives in
\eqref{analytic:eq:free-hypersurfaceA}.
\end{proof}

For a general nonzero $F\in\A_n$, first factor out $t^{\vH(F)}$; then a
complex-linear change of coordinates makes the leading germ $w$-regular,
by choosing a direction on which the first nonzero homogeneous part does not
vanish. Such a change acts coefficientwise and preserves all exponent
supports. These operations reduce any nonunit to the situation of
\cref{analytic:thm:prepA} with $m\ge1$.

\begin{corollary}[Fibres of a prepared hypersurface over $\A_n$]
\label{analytic:cor:hypersurface-fiber}
For $a'\in\m_K^{n-1}$ the fibre algebra of
\eqref{analytic:eq:free-hypersurfaceA} is $K[w]/P(a',w)$. It has length $m$,
all its roots are infinitesimal, and it has $m$ distinct roots exactly when
the discriminant of $P(a',w)$ is nonzero.
\end{corollary}
\begin{proof}
The lower coefficients of $P(a',w)$ lie in $\m_K$: their exponent-zero
coefficient germs vanish at the origin in the $x$ variables, and all
remaining Hahn terms have positive valuation. A root of negative valuation
would make $w^m$ the unique term of least valuation; a root of valuation
zero would reduce to a nonzero root of $w^m$ over $\C$. Both are
impossible. Algebraic closedness of $K$ and the discriminant criterion
finish the proof.
\end{proof}

This is a genuine finite-mapping statement with infinitesimal parameter
values, not merely a one-variable factorization at a single point. It does
not assert that the roots can be globally labelled by single-valued
analytic functions through a branch locus.

\subsection{Division over the common-domain ring $\RR_n$}
\label{analytic:subsec:divR}
The ordinary input here is different, and it is what keeps the domain
fixed. Write $z=(x,w)$ with $x\in\C^{n-1}$, let $V$ be an ordinary polydisk
and $D_R=\{|w|<R\}$, and let $P_0(x,w)$ be monic of degree $d$ in $w$ with
coefficients holomorphic on $V$ and every root of $P_0(x,\cdot)$ in $D_r$
with $r<R$ fixed. Choose $r<\rho<R$ and set, for $h\in\OO(V\times D_R)$,
\begin{equation}\label{analytic:eq:ordinaryR}
 R_0h(x,w)=\frac1{2\pi i}\int_{|\zeta|=\rho}
 \frac{h(x,\zeta)}{P_0(x,\zeta)}\cdot
 \frac{P_0(x,\zeta)-P_0(x,w)}{\zeta-w}\,d\zeta ,
 \qquad
 D_0h=\frac{h-R_0h}{P_0}.
\end{equation}
The quotient involving $P_0$ is a polynomial in $w$ of degree below $d$, so
$R_0h$ has that degree with coefficients holomorphic on $V$; for $|w|<\rho$
the Cauchy formula shows $h-R_0h$ is divisible by $P_0$, and since all zeros
of $P_0$ lie in the smaller disk the quotient extends holomorphically to all
of $V\times D_R$. Thus
\begin{equation}\label{analytic:eq:ordinarydivision}
 h=R_0h+P_0D_0h,\qquad \deg_wR_0h<d,
\end{equation}
uniquely, and \emph{both operators are complex-linear endomorphisms of the
same space} $\OO(V\times D_R)$. That last clause is the resource unavailable
over $\A_n$.

\begin{theorem}[Parameterized Hahn preparation over $\RR_n$]
\label{analytic:thm:prepR}
Let $F\in\RR_n$ be nonzero, normalized as $F=t^\lambda(f_0+E)$ with
$\vH(E)>0$, and suppose $f_0(0,w)$ has a zero of order $d\ge1$ at $w=0$.
Choose an ordinary product neighbourhood on which ordinary preparation gives
$f_0=u_0P_0$ with $u_0$ nowhere zero, $P_0$ monic of degree $d$ and
$P_0(0,w)=w^d$. On a possibly smaller product neighbourhood there are unique
normalized factors
\begin{equation}\label{analytic:eq:prepR}
 F=t^\lambda u_0\,(P_0+A)(1+B),\qquad \deg_wA<d,
\end{equation}
with $A$ and $B$ of strictly positive Hahn support. If $E'=E/u_0$ has
support $E_+$, then $\supp A,\supp B\subseteq E_+^*\setminus\{0\}$. Every
coefficient function of every factor is holomorphic on this one
neighbourhood.
\end{theorem}
\begin{proof}
Ordinary Weierstrass preparation is used once, on $f_0$
\cite{GrauertRemmert}. Shrink the parameter polydisk so that the roots of
$P_0$ stay in a fixed smaller $w$-disk, giving the operators
\eqref{analytic:eq:ordinaryR}. Dividing by $t^\lambda u_0$, the equation to
solve is $E'=A+P_0B+AB$. Proceed by well-founded recursion on
$E_+^*\setminus\{0\}$: at exponent $\nu$ put
\begin{equation}\label{analytic:eq:preprecursion}
 h_\nu=E'_\nu-\!\!\sum_{\substack{\alpha+\beta=\nu\\ \alpha,\beta>0}}\!\!
 A_\alpha B_\beta,\qquad
 A_\nu=R_0h_\nu,\qquad B_\nu=D_0h_\nu .
\end{equation}
The displayed sum is finite by \cref{analytic:lem:neumann} and uses only
smaller exponents, and \eqref{analytic:eq:ordinarydivision} gives the
desired equation at $\nu$. The support lemma proves the assembled factors
exist. Crucially $R_0,D_0$ never change the chosen ordinary domain. For
uniqueness, compare two normalized pairs at the least exponent where either
differs: the product $AB$ cannot contribute a first difference there, since
both factors have positive support, so $\Delta A_\nu+P_0\Delta B_\nu=0$ and
ordinary division uniqueness makes both differences zero. The factor $1+B$
is a unit by positive-support geometric inversion and $u_0$ is an ordinary
unit.
\end{proof}

\begin{theorem}[Parameterized Hahn division over $\RR_n$]
\label{analytic:thm:divisionR}
Let $P=P_0+A$ be the monic factor of \cref{analytic:thm:prepR} and let $G$
be a Hahn-holomorphic datum on its product neighbourhood. There are unique
$Q$ and $R$ on that same neighbourhood with
\begin{equation}\label{analytic:eq:divisionR}
 G=PQ+R,\qquad\deg_wR<d,
\end{equation}
and with $E_A=\supp A\subset\Gamma_{>0}$ both outputs have support in
$\supp G+E_A^*$. They are given by
\begin{equation}\label{analytic:eq:divisionR-formula}
 Q=\sum_{k\ge0}\bigl(-D_0(A\,\cdot)\bigr)^kD_0G,
 \qquad R=R_0(G-AQ).
\end{equation}
For an arbitrary germ $G\in\RR_n$ one first chooses a common smaller
product neighbourhood; no subsequent shrinking is needed.
\end{theorem}
\begin{proof}
Apply $D_0$ to \eqref{analytic:eq:divisionR}: the equation for $Q$ is
$Q+D_0(AQ)=D_0G$, and \cref{analytic:lem:operator} gives
\eqref{analytic:eq:divisionR-formula} with its support bound. Set
$R=R_0(G-AQ)$ and use \eqref{analytic:eq:ordinarydivision} to recover
\eqref{analytic:eq:divisionR}. If $PQ+R=0$ with $\deg_wR<d$, the same
invertible operator gives $Q=0$ and then $R=0$.
\end{proof}

\begin{corollary}[Finite free hypersurface quotient over $\RR_n$]
\label{analytic:cor:freehypersurfaceR}
For $P$ as above,
$\RR_n/(P)\simeq\bigoplus_{j=0}^{d-1}\RR_{n-1}w^j$ as an
$\RR_{n-1}$-module, and for each $a'\in\m_K^{n-1}$ all $d$ roots of
$P(a',w)$ are infinitesimal, counted with multiplicity.
\end{corollary}
\begin{proof}
Division gives the module isomorphism. The lower coefficients of
$P_0(a',w)$ are infinitesimal because their ordinary values at $a'=0$
vanish, and the coefficients of $A(a',w)$ are infinitesimal as well. If a
root $w$ were not infinitesimal, dividing the equation by $w^d$ would give
$1$ plus an infinitesimal equal to zero. Algebraic closedness of $K$
supplies all $d$ roots.
\end{proof}

\begin{remark}[Why both division theorems are kept]
\label{analytic:rem:two-divisions}
\Cref{analytic:thm:divisionA} applies to coefficient families with radii
tending to zero, which \cref{analytic:thm:divisionR} cannot reach; the
element \eqref{analytic:eq:no-radius} is such a family.
\Cref{analytic:thm:divisionR} delivers its quotient and remainder on
\emph{one named ordinary domain}, which \cref{analytic:thm:divisionA}
cannot deliver, and that fixed domain is what later permits Cartan's
theorem~B (\cref{analytic:lem:ordinarykoszul}) and a fixed ordinary residue
cycle (\cref{analytic:thm:residueduality}). Neither theorem implies the
other.
\end{remark}

\section{Noetherian geometry and the Nullstellensatz over $\A_n$}
\label{analytic:sec:structure-A}

Everything in this section is a statement about the radius-free ring
$\A_n=\C\{z\}((t^\Gamma))$ of \cref{analytic:def:ringA}. The corresponding
statements for $\RR_n$ are in \cref{analytic:sec:structure-R} and are not
consequences of these.

\subsection{Finite generation}
\begin{theorem}[Noetherianity of $\A_n$]\label{analytic:thm:noetherianA}
For every $n\ge0$ the ring $\A_n$ is Noetherian.
\end{theorem}
\begin{proof}
Induct on $n$, starting from the field $\A_0=K$. Let $I\subseteq\A_n$ be a
nonzero proper ideal and choose $0\ne F\in I$. Normalize and change
complex-linear coordinates so that \cref{analytic:thm:prepA} applies with
degree $m\ge1$. Then $\A_n/(F)$ is finite free over $\A_{n-1}$, hence a
Noetherian $\A_{n-1}$-module by the induction hypothesis, so its submodule
$I/(F)$ is finitely generated over $\A_{n-1}$. Lifting those generators and
adjoining $F$ gives finitely many ideal generators of $I$. The zero and unit
ideals need no argument.
\end{proof}

\subsection{Every maximal ideal of $\A_n$ is an evaluation ideal}
\begin{theorem}[Weak Nullstellensatz over $\A_n$]
\label{analytic:thm:weak-nullA}
Every maximal ideal of $\A_n$ is $\m_a$ for a unique $a\in\m_K^n$, and its
residue field is $K$.
\end{theorem}
\begin{proof}
The case $n=0$ is clear. Assume the statement in dimension $n-1$ and let $M$
be maximal in $\A_n$. Since $\A_n$ is not a field for $n\ge1$, $M$ contains
a nonzero element; make it regular as above, obtaining the finite integral
algebra $\A_n/(F)\simeq\A_{n-1}[w]/(P)$. The contraction of $M/(F)$ to
$\A_{n-1}$ is maximal by integrality, hence by induction is evaluation at
some $a'\in\m_K^{n-1}$ with residue field $K$; therefore $\A_n/M$ is a
finite field extension of $K$ and equals $K$. Let $a_n$ be the image of $w$
in that field; it satisfies $P(a',a_n)=0$, so
\cref{analytic:cor:hypersurface-fiber} gives $a_n\in\m_K$. The coordinate
images thus form an infinitesimal tuple $a$; the ideal $(z-a)$ is contained
in $M$ and is maximal by \cref{analytic:prop:kernels}, so $M=(z-a)$.
Reversing the complex-linear coordinate change keeps the coordinates
infinitesimal, and evaluating the coordinate functions gives uniqueness.

This proof never assumes that an abstract field-valued homomorphism
commutes with an infinite Hahn sum. The finite integral quotient identifies
the coordinate images, and the independently proved kernel formula
\eqref{analytic:eq:kernel} identifies the ideal.
\end{proof}

\subsection{The strong theorem}
For $X\subseteq\m_K^n$ put $I(X)=\{F\in\A_n:F(a)=0\ \forall a\in X\}$, and
for an ideal $I$ let $Z(I)\subseteq\m_K^n$ be its common zero set.

\begin{theorem}[Jacobson property and strong Nullstellensatz over $\A_n$]
\label{analytic:thm:strong-nullA}
The ring $\A_n$ is Jacobson, and for every ideal $I\subseteq\A_n$
\begin{equation}\label{analytic:eq:strong-nullA}
 I(Z(I))=\sqrt I .
\end{equation}
\end{theorem}
\begin{proof}
We prove the Jacobson assertion by induction on $n$, the base being a
field. Let $\pp\ne0$ be prime, choose $0\ne F\in\pp$ and prepare it after a
coordinate change; then $\A_n/\pp$ is finite over
$\A_{n-1}/(\pp\cap\A_{n-1})$, which is Jacobson by induction and passage to
a quotient, and a finite algebra over a Jacobson ring is Jacobson
\cite[\S10.35]{Stacks}. Hence $\pp$ is the intersection of the maximal
ideals containing it.

The zero prime requires a separate argument; the finite-quotient reduction
does not apply to it. By \cref{analytic:prop:faithful} the intersection of
all evaluation ideals $\m_a$ is zero. Thus every prime is an intersection of
maximal ideals, which is the Jacobson property, and by
\cref{analytic:thm:weak-nullA} those maximal ideals are precisely the
evaluation ideals. Intersecting the maximal ideals containing $I$ gives
\eqref{analytic:eq:strong-nullA}.
\end{proof}

A family of Hahn-germ equations with no common infinitesimal zero therefore
generates the unit ideal, and because $\A_n$ is Noetherian the assertion has
a finite witness: finitely many $F_j$ and $G_j\in\A_n$ with
$\sum_jG_jF_j=1$. The theorem gives existence, not a uniform bound on the
complexity of finding such a witness from arbitrary support presentations.

\subsection{Regularity, completion, and two warnings}
\begin{theorem}[Local geometry of $\A_n$]\label{analytic:thm:regularA}
The ring $\A_n$ has dimension $n$ and is regular. For every $a\in\m_K^n$,
\begin{equation}\label{analytic:eq:completionA}
 \widehat{(\A_n)_{\m_a}}\simeq K[[X_1,\ldots,X_n]],\qquad X_i=z_i-a_i .
\end{equation}
\end{theorem}
\begin{proof}
Translate to $a=0$. For $0\le r\le n$, coefficientwise restriction gives
$\A_n/(z_1,\ldots,z_r)\simeq\A_{n-r}$, a domain, so the chain of coordinate
primes remains strict after localizing at $\m_0$ and that local ring has
dimension at least $n$. Its maximal ideal is generated by the $n$ coordinate
functions, so the standard Noetherian local dimension bound gives dimension
at most $n$, and the finite-jet calculation \eqref{analytic:eq:jets} shows
the cotangent space has dimension $n$; hence the local ring is regular. By
\cref{analytic:thm:weak-nullA} every maximal localization is of this form,
so $\A_n$ has dimension $n$ and all localizations at primes are regular,
regularity being preserved under localization. Finally an element outside
$\m_a$ is already invertible modulo each $\m_a^N$, since its constant term
there is nonzero, so localization does not change the finite-jet quotients
\eqref{analytic:eq:jets}; taking the inverse limit gives
\eqref{analytic:eq:completionA}.
\end{proof}

\begin{remark}[$\A_n$ is not local]\label{analytic:rem:nonlocalA}
The ring $\A_n$ simultaneously records every infinitesimal point, so it has
one maximal ideal for each $a\in\m_K^n$; its ordinary coefficient ring $R_n$
is local by contrast. This is exactly why an infinitesimal Nullstellensatz
can coexist with nontrivial local analytic geometry, and it is the formal
content of the distinction between a stalk on the whole monad and a fine
germ at one surcomplex point. Nothing below identifies the two.
\end{remark}

\begin{remark}[The nonzero-value-group hypothesis is essential]
\label{analytic:rem:trivial-group}
If $\Gamma=\{0\}$ then $\m_K^n=\{0\}$ and $\A_n=R_n$. For $n>0$ the
vanishing ideal of that single point is $(z)$, whereas $\sqrt{(0)}=(0)$;
thus \eqref{analytic:eq:strong-nullA} fails at $I=(0)$. The hypothesis is
used essentially in \cref{analytic:prop:faithful}. It is not an aesthetic
restriction.
\end{remark}

\section{Noetherian geometry, normalization and finite maps over $\RR_n$}
\label{analytic:sec:structure-R}

Everything in this section is a statement about the common-domain ring
$\RR_n$ of \cref{analytic:def:ringR}. The statements parallel those of
\cref{analytic:sec:structure-A} but are proved by a different route, and
the normalization and finite-map results of
\cref{analytic:sec:finite} have no counterpart there.

\begin{theorem}[Noetherian monad ring $\RR_n$]
\label{analytic:thm:noetherianR}
The ring $\RR_n$ is a Noetherian domain of Krull dimension $n$.
\end{theorem}
\begin{proof}
The domain assertion follows from multiplication of leading coefficient
germs. For Noetherianity induct on $n$ from $\RR_0=K$. Let $I\subset\RR_n$
be nonzero and choose $0\ne F\in I$. An ordinary invertible linear change of
coordinates makes its leading ordinary coefficient regular in the last
variable, by choosing a direction where its first nonzero homogeneous
Taylor polynomial does not vanish. If that coefficient has order zero then
$F$ is a unit after shrinking the ordinary neighbourhood and $I=\RR_n$.
Otherwise \cref{analytic:thm:prepR} replaces $F$, up to a unit, by a monic
$P$ of positive degree in the last variable, and by
\cref{analytic:cor:freehypersurfaceR} the quotient $\RR_n/(F)$ is finite
free over $\RR_{n-1}$; as that ring is Noetherian, $I/(F)$ is finitely
generated, and lifting generators and adjoining $F$ generates $I$.

For the dimension, the coordinate ideals
$(0)\subsetneq(z_1)\subsetneq\cdots\subsetneq(z_1,\ldots,z_n)$ are prime by
\eqref{analytic:eq:coordinatequotient}, giving the lower bound $n$. For the
upper bound, prepend $(0)$ to any prime chain if necessary and choose a
nonzero element of its first nonzero prime; preparation makes the quotient
by that element finite integral over $\RR_{n-1}$, and integral extensions
preserve dimension, so the remaining chain has length at most $n-1$. The
elementary integral-extension facts used here --- incomparability, going up
and lying over --- are the standard results recorded in \cite{Stacks}.
\end{proof}

\begin{proposition}[Noether normalization over $\RR_n$]
\label{analytic:prop:normalization}
Let $\pp\subset\RR_n$ be prime and $d=\dim(\RR_n/\pp)$. After an ordinary
linear coordinate change, the first $d$ coordinates induce a finite
injective homomorphism
\begin{equation}\label{analytic:eq:normalization}
 \RR_d\hookrightarrow A=\RR_n/\pp .
\end{equation}
\end{proposition}
\begin{proof}
Induct on $n$. If $\pp=0$ take $d=n$. Otherwise prepare a nonzero element of
$\pp$ in the last variable; then $A$ is finite over $\RR_{n-1}/\mathfrak q$
with $\mathfrak q$ the contraction of $\pp$, and by induction that domain is
finite over an injective coordinate copy of some $\RR_d$. Finiteness is
transitive. All coordinate changes in the induction act ordinarily on the
remaining coordinates, so their composition is an ordinary linear change on
$\C^n$; dimension is preserved in each finite integral extension, which
identifies the final number of coordinates with $\dim A$.
\end{proof}

This proposition has no counterpart over $\A_n$ in this report, and it is
the input to both the strong Nullstellensatz proof below and the finite-map
theorem of \cref{analytic:sec:finite}.

\begin{theorem}[Weak monad Nullstellensatz over $\RR_n$]
\label{analytic:thm:weaknullR}
Every maximal ideal of $\RR_n$ is uniquely
$\m_a=(z_1-a_1,\ldots,z_n-a_n)$ with $a\in\m_K^n$; its residue field is $K$
and its residue map is Hahn evaluation at $a$.
\end{theorem}
\begin{proof}
Induct on $n$, the case $n=0$ being immediate. For $n>0$, $\RR_n$ is not a
field since evaluation at zero shows $z_1$ is a nonunit, so a maximal ideal
$\m$ is nonzero. Prepare an element of $\m$ and use the resulting finite
extension of $\RR_{n-1}$; the contraction $\mathfrak n$ is maximal, since a
subring over which a field is integral is a field, and by induction
$\RR_{n-1}/\mathfrak n=K$. Hence $\RR_n/\m$ is a finite extension of the
algebraically closed $K$, so equals $K$.

Let $a_j\in K$ be the image of $z_j$. If $v(a_j)<0$ then $z_j-a_j$ has a
nonzero constant as leading Hahn coefficient after normalization and is a
unit. If $v(a_j)=0$ its leading coefficient is $z_j-c$ with $c\in\C^\times$,
a unit on a sufficiently small ordinary neighbourhood of zero. Both cases
contradict $z_j-a_j\in\m$, so every $a_j$ is infinitesimal or zero. By
\eqref{analytic:eq:coordinatequotient} the ideal generated by these
coordinate differences already has quotient $K$ and is precisely the
evaluation kernel, so it equals $\m$. The coordinate images give
uniqueness.
\end{proof}

\begin{theorem}[Strong monad Nullstellensatz over $\RR_n$]
\label{analytic:thm:strongnullR}
For an ideal $I\subset\RR_n$ put
$V(I)=\{a\in\m_K^n:F(a)=0\ \text{for all}\ F\in I\}$. Then
\begin{equation}\label{analytic:eq:strongnullR}
 I(V(I))=\sqrt I ,
\end{equation}
equivalently $\RR_n$ is Jacobson and all its closed points are its
infinitesimal evaluation points.
\end{theorem}
\begin{proof}
We first show a prime $\pp$ is the intersection of the maximal ideals
containing it. Let $A=\RR_n/\pp$ and $0\ne g\in A$, and normalize $A$ over
$B=\RR_d$ as in \cref{analytic:prop:normalization}. The finite-dimensional
domain $A\otimes_B\Frac(B)$ is a field, so $g^{-1}$ lies in it; clearing a
denominator gives $gh=b$ with $h\in A$ and $0\ne b\in B$. By
\cref{analytic:lem:density} choose an evaluation maximal ideal
$\mathfrak n$ of $B$ not containing $b$; lying over for the integral
extension $B\subset A$ gives a maximal ideal $M$ of $A$ over
$\mathfrak n$, and $gh=b$ implies $g\notin M$. Pulling $M$ back to $\RR_n$
and applying \cref{analytic:thm:weaknullR} produces a point of $V(\pp)$
where any representative of $g$ is nonzero.

If $F\notin\sqrt I$, choose a prime $\pp\supseteq I$ not containing $F$; the
preceding argument finds an evaluation maximal ideal containing $\pp$ but
not $F$, so $F$ does not vanish on all of $V(I)$. The reverse inclusion in
\eqref{analytic:eq:strongnullR} is immediate in the field $K$.
\end{proof}

\begin{remark}[The two Jacobson proofs are genuinely different]
\label{analytic:rem:two-jacobson}
\Cref{analytic:thm:strong-nullA} obtains the Jacobson property of $\A_n$ by
induction through finite algebras over Jacobson rings, and then handles the
zero prime separately by \cref{analytic:prop:faithful}, because the
finite-quotient reduction does not apply to it.
\Cref{analytic:thm:strongnullR} instead runs Noether normalization, inverts
$g$ in the generic fibre, and uses lying over together with
\cref{analytic:lem:density} to produce the required point; the zero prime
needs no separate treatment because normalization applies to it. Neither
argument is a specialization of the other, and neither theorem may be
transferred between the rings.
\end{remark}

\begin{proposition}[Local coordinates and completion over $\RR_n$]
\label{analytic:prop:completionR}
For $a\in\m_K^n$ the local ring $(\RR_n)_{\m_a}$ is regular of dimension
$n$, and
\begin{equation}\label{analytic:eq:completionR}
 \widehat{(\RR_n)_{\m_a}}\simeq K[[h_1,\ldots,h_n]],\qquad h_j=z_j-a_j .
\end{equation}
The image of $F$ in this completion is its formal Taylor series at $a$.
\end{proposition}
\begin{proof}
Translate $a$ to zero. The maximal ideal is generated by the $n$ coordinate
functions and the coordinate prime chain survives localization, so the local
dimension is at least $n$, hence exactly $n$ by
\cref{analytic:thm:noetherianR}; a Noetherian local ring of dimension $n$
whose maximal ideal has $n$ generators is regular. For every $r\ge1$,
ordinary Taylor division applied coefficientwise gives a unique
decomposition into the polynomial of total degree below $r$ and a member of
$(z_1,\ldots,z_n)^r$, whence
$\RR_n/\m_0^r\simeq K[z]/(z)^r$. Localizing does not change this quotient,
whose radical is already maximal, and the inverse limit over $r$ gives
\eqref{analytic:eq:completionR}. Translation and
\cref{analytic:prop:evaluationR} identify the coefficients with the Taylor
coefficients at $a$.
\end{proof}

Consequently, for an isolated zero of $F_1,\ldots,F_n$ at $a$ the local
intersection multiplicity is unambiguously
\begin{equation}\label{analytic:eq:localmultiplicity}
 e_a=\dim_K\frac{K[[h_1,\ldots,h_n]]}{(F_1(a+h),\ldots,F_n(a+h))}
\end{equation}
when finite, and it equals the length of the corresponding local quotient of
$\RR_n$, since completion preserves finite-length modules. The same formula
holds over $\A_n$ by \cref{analytic:thm:regularA}; it is one formula
instantiated twice, not one theorem transferred.

\begin{remark}[Neither ring is local]\label{analytic:rem:nonlocalR}
\Cref{analytic:thm:strongnullR} is not the ordinary complex local
Nullstellensatz restated with a different scalar field: $\RR_n$ has many
maximal ideals, whereas $\C\{z\}$ is local. The same remark applies to
$\A_n$ (\cref{analytic:rem:nonlocalA}). This is the sense in which these
rings record an entire infinitesimal monad rather than a single fine germ.
\end{remark}

\section{Finite maps of monad zero sets over $\RR_n$}
\label{analytic:sec:finite}

\begin{theorem}[Finite projection and generic fibres over $\RR_n$]
\label{analytic:thm:finitemap}
Let $\pp\subset\RR_n$ be prime and let $A=\RR_n/\pp$ have dimension $d$.
After an ordinary linear coordinate change the coordinate projection
\[
 \pi:V(\pp)\longrightarrow\m_K^d
\]
is surjective with finite fibres, induced by a finite injective map
$B=\RR_d\hookrightarrow A$. If $A$ is generated by $N$ elements as a
$B$-module, every fibre has scheme length at most $N$. Let
$r=[\Frac(A):\Frac(B)]$. There is a nonzero $s\in B$ such that every
$b\in\m_K^d$ with $s(b)\ne0$ has exactly $r$ distinct preimages, each
reduced; such $b$ exist.
\end{theorem}
\begin{proof}
Use \cref{analytic:prop:normalization}. Lying over makes the map on maximal
ideals surjective, and by \cref{analytic:thm:weaknullR} these are exactly
the evaluation points on both sides, which gives surjectivity on monads. For
an evaluation maximal ideal $\mathfrak n_b\subset B$ the fibre algebra
$A/\mathfrak n_bA$ is a nonzero finite-dimensional $K$-algebra spanned by
the $N$ module generators, giving the length bound and finiteness.

Set $L=\Frac(A)$, $L_0=\Frac(B)$. The extension $L/L_0$ is finite and
separable since the characteristic is zero, so by the primitive element
theorem and clearing denominators there is a primitive element $\theta$ in a
localization of $A$. The finitely many module generators of $A$ are
$L_0$-linear combinations of $1,\theta,\ldots,\theta^{r-1}$. Localize at one
nonzero $s\in B$ so that these expressions, the coefficients of the monic
minimal polynomial $q(T)$ of $\theta$, and $\theta$ itself are defined over
$B_s$, and enlarge $s$ to make the nonzero discriminant of $q$ invertible.
Then
\[
 A_s\simeq B_s[T]/(q),\qquad\deg q=r,\qquad \disc(q)\in B_s^\times :
\]
the expressions for the generators give surjectivity, and monic division
together with independence of $1,\theta,\ldots,\theta^{r-1}$ over $L_0$
gives injectivity. Each such fibre is therefore the algebra of a degree-$r$
polynomial with nonzero discriminant over the algebraically closed $K$, so
it has $r$ distinct points. Finally \cref{analytic:lem:density} supplies a
$b$ avoiding $s=0$.
\end{proof}

The word \emph{finite} here has its algebraic meaning throughout: finite
module and finite fibres. No claim of properness, compactness, or classical
path lifting in the full surreal fine topology is included. Two senses of
finiteness occur in this report and both are purely algebraic and
free-of-finite-rank, never metric: the \emph{finite projection} of
\cref{analytic:thm:prepA} and \cref{analytic:cor:freehypersurfaceR}, in
which $\A_n/(F)$ respectively $\RR_n/(P)$ is free of rank $m$ over the ring
in one fewer variable with basis $1,w,\ldots,w^{m-1}$; and the \emph{finite
fibre} of \cref{analytic:cor:hypersurface-fiber} and
\cref{analytic:thm:finitemap}, an algebra of length $m$ over $K$ all of
whose roots are infinitesimal.

\begin{remark}[What $\A_n$ is not given here]
\label{analytic:rem:no-normalization-A}
We do not assert a Noether normalization or a finite-projection theorem for
$\A_n$. The normalization proof above uses preparation with holomorphic
parameters on one fixed neighbourhood, and we have not verified that the
germ-level preparation of \cref{analytic:thm:prepA} supports the same
induction. The hypersurface case \cref{analytic:cor:hypersurface-fiber} is
available over $\A_n$, and nothing more is claimed.
\end{remark}

\section{Isolated complete intersections: the classical ingredients}
\label{analytic:sec:ordinary}

The deformation package --- division with the $S+E^*$ certificate,
conservation of $\mu$, characters as actual infinitesimal evaluations, the
perfect residue pairing and the trace--Jacobian identity --- is common
property of this subject and is stated at its strongest hypotheses in the
companion report on finite deformations \cite{FiniteDeformations}. What is
distinctive here, and what this report contributes, is that the package
survives the \emph{removal of the common-radius hypothesis}: it holds over
the strictly larger ring $\A_n$, which by \cref{analytic:ex:no-radius}
contains elements no common-domain theorem can reach. We therefore state
the theorem over each ring, give the two proofs, and refer to
\cite{FiniteDeformations} for the parts of the package that are neither
strengthened nor re-proved here.

Throughout \crefrange{analytic:sec:ordinary}{analytic:sec:extension}, let
$f=(f_1,\ldots,f_n)\in R_n^n$ vanish at the origin with
\begin{equation}\label{analytic:eq:mu}
 0<\mu=\dim_\C R_n/(f_1,\ldots,f_n)<\infty .
\end{equation}
This is the algebraic form of an isolated common zero. Since $R_n$ is
regular local and the ideal is generated by $n$ elements with
zero-dimensional quotient, it is a complete intersection. Choose monomials
$b_1,\ldots,b_\mu$ whose classes form a $\C$-basis of $B_0=R_n/(f)$; such a
choice exists because some power of $(z)$ lies in $(f)$, so finitely many
monomials span the quotient and a subset is a basis.

\subsection{A linear division splitting}\label{analytic:subsec:splitting}
Let $i:\C^\mu\to R_n$ send coordinate vectors to the chosen monomials and
let $p:R_n\to\C^\mu$ be the coordinate map of the quotient in this basis.
There is a complex-linear $h:R_n\to R_n^n$ with
\begin{equation}\label{analytic:eq:splitting}
 g=i\,p(g)+\sum_{j=1}^nf_jh_j(g),\qquad pi=\id,\qquad hi=0 :
\end{equation}
take a complex-linear right inverse of the surjection $R_n^n\to(f)$ and
apply it to $g-ip(g)$. This is a choice of linear splitting, not a claim of
uniqueness of the quotients $h_j(g)$; one may instead use a fixed
constructive local division procedure when an effective presentation is
available. No ordinary norm bound on this splitting is assumed. It maps
analytic germs to analytic germs, which is exactly what coefficientwise Hahn
extension requires, and this freedom is what removes the need to choose one
analytic neighbourhood for infinitely many input coefficients.

\subsection{Local residues and duality, imported}
\label{analytic:subsec:classical-residue}
We explicitly import the following classical facts; they are not novelty
claims. For an isolated tuple $g=(g_1,\ldots,g_n)$ in $R_n$ the ordinary
local Grothendieck residue is the complex-linear functional
\begin{equation}\label{analytic:eq:classical-residue}
 \res_g(q)=\frac1{(2\pi i)^n}\int_{\{|g_j|=\epsilon_j\}}
 \frac{q(z)\,dz_1\wedge\cdots\wedge dz_n}{g_1\cdots g_n},
\end{equation}
with sufficiently small regular radii and the orientation given by
increasing arguments of the $g_j$. This is an \emph{ordinary}
complex-valued contour integral over an ordinary real cycle in $\C^n$;
below it is extended coefficientwise to Hahn numerators, and this report
never constructs a surcomplex contour. For background and the equivalent
algebraic treatment see \cite{Tsikh,SS,CDS,GH}.

The residue annihilates $(g)$ and induces a nondegenerate bilinear pairing
on $R_n/(g)$. For coordinate powers,
\begin{equation}\label{analytic:eq:coordinate-residue}
 \res_{(z_1^{d_1},\ldots,z_n^{d_n})}(q)=[z_1^{d_1-1}\cdots z_n^{d_n-1}]\,q .
\end{equation}
For $k\in\N^n$ write $f^{k+1}=(f_1^{k_1+1},\ldots,f_n^{k_n+1})$;
cancellation of a denominator power gives
\begin{equation}\label{analytic:eq:pole-cancel}
 \res_{f^{k+1}}(f_jq)=
 \begin{cases}
 0,&k_j=0,\\
 \res_{f^{k+1-e_j}}(q),&k_j>0,
 \end{cases}
\end{equation}
the denominator tuple in the second line still having every exponent
positive. These are special cases of the residue transformation rule. We
also use the ordinary finite-family trace identity: in a sufficiently small
analytic deformation $g_u$ of $f$ by finitely many ordinary parameters $u$,
the cluster originating at zero has a finite free algebra of rank $\mu$ over
the parameter-germ ring and
\begin{equation}\label{analytic:eq:ordinary-trace}
 \res^{\mathrm{cluster}}_{g_u}(qJ_{g_u})=\Tr(M_q),
\end{equation}
the cluster residue being the sum over all points in a fixed small
representative neighbourhood, multiple points included. Its algebraic form
is the trace--residue theorem of \cite{SS}; deformation and computational
formulations appear in \cite{KW,CDS}. A full justification of the
finite-family input and its exact role is in
\cref{analytic:app:classical-family}.

The classical facts above are used on \emph{one finite family at a time}.
We never assert that all coefficient germs of a Hahn perturbation admit a
common contour or a common analytic representative. Over $\RR_n$ a common
representative does exist, by definition of the ring, and
\cref{analytic:thm:residueduality} exploits it; over $\A_n$ it does not, and
\cref{analytic:lem:finite-dependency} is the substitute.

\subsection{The perturbation residue and its normalization}
\label{analytic:subsec:normalization}
Keep \eqref{analytic:eq:mu} and put
\begin{equation}\label{analytic:eq:perturbed}
 F_j=f_j+E_j,\qquad \vH(E_j)>0,\qquad
 E=\bigcup_j\supp E_j\subset\Gamma_{>0},
\end{equation}
with $E_j$ in $\A_n$ or in $\RR_n$ as the context specifies. Extending each
ordinary residue functional coefficientwise to Hahn numerators, define for
$H$ in the ambient ring
\begin{equation}\label{analytic:eq:hahn-residue}
 \boxed{\quad
 \Res_F(H)=\sum_{k\in\N^n}(-1)^{|k|}
 \res_{f^{k+1}}\bigl(H\,E_1^{k_1}\cdots E_n^{k_n}\bigr).\quad}
\end{equation}
The sign $(-1)^{|k|}$, together with the factor $1/(2\pi i)^n$ on the
contour side of \eqref{analytic:eq:classical-residue}, is the
normalization, and it is chosen so that $\Res_F(J_F)=\mu$ and, at simple
zeros, $\Res_F(H)=\sum_aH(a)/J_F(a)$.

We state emphatically, once: \eqref{analytic:eq:hahn-residue} is
\emph{motivated} by expanding each reciprocal $(f_j+E_j)^{-1}$, but the
displayed series, not a putative surcomplex contour integral, \emph{is} the
definition. Supplying an actual surcomplex contour that represents it is a
separate task, addressed in the companion report on contours
\cite{Contours}. The functional $\Lambda_F$ of
\cref{analytic:subsec:contour-residue} is not a counterexample to this: it
expands an \emph{ordinary} cycle coefficientwise, and it is defined only
over $\RR_n$, where such a cycle exists.

Three senses of multiplicity are kept apart throughout and should not be
conflated:
\begin{enumerate}
\item the \emph{global colength} $\mu=\dim_\C R_n/(f)=\dim_KB_F$, the total
length, conserved;
\item the \emph{local algebraic multiplicity} $e_a=\dim_KB_a$ in the Artin
decomposition $B_F\simeq\prod_{a\in Z_F}B_a$, with $e_a\ge1$ and
$\sum_ae_a=\mu$;
\item the same $e_a$ computed \emph{formally at the actual displaced zero}
$a$ rather than at its standard part, by
\eqref{analytic:eq:localmultiplicity}.
\end{enumerate}

\section{The deformation theorem over each ring}\label{analytic:sec:deform}

\subsection{The radius-free form, over $\A_n$}
\label{analytic:subsec:deformA}
\begin{lemma}[Well-definedness and annihilation over $\A_n$]
\label{analytic:lem:residue-annihilate}
For $E_j\in\A_n$, the sum \eqref{analytic:eq:hahn-residue} is strongly
summable and defines a $K$-linear functional $\A_n\to K$. If $S=\supp H$
then $\supp\Res_F(H)\subseteq S+E^*$. For every $j$ and $H$,
$\Res_F(F_jH)=0$.
\end{lemma}
\begin{proof}
The ordinary residue operators do not enlarge Hahn support. A term of
multi-index $k$ is supported on sums of a source exponent from $S$ and
$|k|$ positive perturbation exponents, so \cref{analytic:lem:neumann}
proves both well-ordering and finiteness of the full family at each
exponent, even though the residue functional itself varies with $k$; this
also gives $K$-linearity by the permitted finite rearrangements.

For annihilation, split $F_jH=f_jH+E_jH$. In the first part, terms with
$k_j=0$ vanish by \eqref{analytic:eq:pole-cancel}; in the remaining terms
cancel one power of $f_j$ and replace $k$ by $k-e_j$. The result is the
negative of the second part of the sum, the reindexing being legitimate
coefficientwise by the finiteness already proved. Hence $\Res_F(F_jH)=0$.
\end{proof}

Extend $i,p,h$ from \eqref{analytic:eq:splitting} coefficientwise and define
the positive-support operator
\begin{equation}\label{analytic:eq:T}
 T:G\longmapsto\sum_{j=1}^nE_j\,h_j(G).
\end{equation}
For $G\in\A_n$ set
\begin{equation}\label{analytic:eq:normal-form}
 L_G=(\id+T)^{-1}G,\qquad
 \NF_F(G)=p(L_G)\in K^\mu,\qquad
 Q_j(G)=h_j(L_G),
\end{equation}
the inverse being the strong Neumann sum of
\cref{analytic:lem:operator}.

\begin{lemma}[Spanning with a support certificate, over $\A_n$]
\label{analytic:lem:spanning}
For every $G\in\A_n$,
\begin{equation}\label{analytic:eq:normal-decomposition}
 G=\sum_{r=1}^\mu\NF_F(G)_r\,b_r+\sum_{j=1}^nF_jQ_j(G),
\end{equation}
and the supports of all displayed coordinates and quotients lie in
$\supp G+E^*$.
\end{lemma}
\begin{proof}
The support claim follows from \cref{analytic:lem:operator} and the
coefficientwise action of $p,h$. Since $G=L_G+TL_G$, applying the ordinary
splitting \eqref{analytic:eq:splitting} to $L_G$ gives
$G=ip(L_G)+\sum_jf_jh_j(L_G)+\sum_jE_jh_j(L_G)$, which is
\eqref{analytic:eq:normal-decomposition}.
\end{proof}

This proves that the $b_r$ \emph{span} the perturbed quotient. Spanning
alone is not conservation of multiplicity: the dimension might still have
dropped. The residue pairing is what prevents that drop.

\begin{theorem}[Finite complete-intersection deformation over $\A_n$:
the radius-free strengthening]\label{analytic:thm:finite-deformA}
Assume \eqref{analytic:eq:mu} and \eqref{analytic:eq:perturbed} with
$E_j\in\A_n$. Then the classes of $b_1,\ldots,b_\mu$ form a $K$-basis of
\[
 B_F=\A_n/(F_1,\ldots,F_n);
\]
in particular $\dim_KB_F=\mu$. The functional $\Res_F$ descends to $B_F$ and
the pairing
\begin{equation}\label{analytic:eq:pairing}
 B_F\times B_F\longrightarrow K,\qquad (u,v)\longmapsto\Res_F(uv)
\end{equation}
is perfect. For the fixed basis, \eqref{analytic:eq:normal-form} is the
unique normal-coordinate map, independently of the auxiliary splitting $h$.
\end{theorem}
\begin{proof}
By \cref{analytic:lem:residue-annihilate} the residue descends to the
quotient. Its Gram matrix in the spanning family is
$\mathcal G_F=(\Res_F(b_rb_s))_{r,s=1}^\mu$. The $k=0$ term of
\eqref{analytic:eq:hahn-residue} is the ordinary Gram matrix
$\mathcal G_0$ of the residue pairing on $B_0$, and all other terms have
strictly positive support; classical local duality gives
$\det\mathcal G_0\ne0$. Hence
\[
 \mathcal G_F=\mathcal G_0+(\text{a positive-support matrix})
\]
is invertible over $K$. If $\sum_rc_rb_r$ lies in $(F)$, pairing with each
$b_s$ gives $\mathcal G_Fc=0$, so all $c_r$ vanish; thus the spanning family
of \cref{analytic:lem:spanning} is a basis and
\eqref{analytic:eq:pairing} is perfect. Any decomposition modulo $(F)$ has
the same coordinates in a fixed basis, which gives independence of $h$. The
quotients $Q_j$ themselves need not be unique for $n>1$.
\end{proof}

The proof has a useful division of labour. A support-controlled Neumann
operator prevents the dimension from increasing; a perturbed nondegenerate
residue matrix prevents it from decreasing. Neither step assumes a simple
zero, a triangular system, a polynomial perturbation, or a common radius of
convergence for the perturbation coefficients. It is the last of these that
is the point of this section: \cref{analytic:thm:finite-deformA} applies to
perturbations such as the coefficient family \eqref{analytic:eq:no-radius},
which lies in $\A_1$ and in no $\HH_\Gamma(U)$.

\subsection{Multiplication matrices and residue reconstruction over $\A_n$}
\label{analytic:subsec:matricesA}
For $H\in\A_n$ define
\begin{equation}\label{analytic:eq:multiplication-matrix}
 M_H=\bigl(\NF_F(Hb_s)_r\bigr)_{r,s=1}^\mu ,
\end{equation}
which is multiplication by the class of $H$ in $B_F$; hence
$M_{H_1H_2}=M_{H_1}M_{H_2}$ and $[M_{z_i},M_{z_j}]=0$. For ordinary basis
elements and coordinate functions the entries have support in $E^*$ and
reduce at exponent zero to their ordinary counterparts. This gives explicit
finite-dimensional data for a deformation whose coefficient support may be
infinite and of arbitrary valuation rank.

The normal form is also reconstructible from residues. With
$r_F(H)=(\Res_F(Hb_s))_{s=1}^\mu$ and the column convention above,
\begin{equation}\label{analytic:eq:residue-reconstruction}
 \NF_F(H)=\mathcal G_F^{-1}r_F(H),
\end{equation}
the symmetry of $\mathcal G_F$ removing any transpose ambiguity. Since
$\mathcal G_0$ is invertible, geometric matrix inversion shows
$\mathcal G_F^{-1}$ has support in $E^*$, so
\eqref{analytic:eq:residue-reconstruction} is a second construction with
exactly the same support bound.

\subsection{The common-domain form, over $\RR_n$: a Koszul conjugacy}
\label{analytic:subsec:deformR}
Over $\RR_n$ a fixed ordinary domain is available, and with it a genuinely
different proof which yields something the $\A_n$ argument does not: the
vanishing of \emph{all} higher homology of the Hahn Koszul complex. We
record it.

Let $(C_\bullet,d)$ be a complex of complex vector spaces concentrated in
nonnegative degrees with homology only in degree zero, and let
$V=H_0(C_\bullet)$ be finite-dimensional. Choose linear maps
$i:V\to C_0$, $p:C_0\to V$ extended by zero in other degrees, and
$h:C_k\to C_{k+1}$, with
\begin{equation}\label{analytic:eq:contraction}
 pi=1,\quad di=0,\quad pd=0,\quad dh+hd=1-ip,\quad
 h^2=0,\quad hi=0,\quad ph=0 .
\end{equation}
Such data always exist: split each space into boundaries and a complement
mapping isomorphically onto the boundaries one degree lower, in degree zero
also splitting off the selected representatives $i(V)$, and define $h$ to
invert $d$ on boundaries and to vanish on the chosen complements and on
$i(V)$. The construction can enforce any specified representatives of a
basis of $V$. These are vector-space splittings: no boundedness or
continuity of them is asserted. Their outputs nevertheless belong to the
same specified space of holomorphic functions, which is the whole point.

\begin{theorem}[Positive Hahn conjugacy]\label{analytic:thm:conjugacy}
Extend the data \eqref{analytic:eq:contraction} coefficientwise to Hahn
families. Let $\delta=\sum_{e\in E}t^e\delta_e$ be a degree-$(-1)$ operator
with $E\subset\Gamma_{>0}$ well ordered, and suppose $(d+\delta)^2=0$. Put
$D=d+\delta$ and
\begin{equation}\label{analytic:eq:perturbedmaps}
 U=1+\delta h,\qquad p_F=pU^{-1},\qquad h_F=hU^{-1}.
\end{equation}
Then
\begin{equation}\label{analytic:eq:chainconjugacy}
 D=UdU^{-1},\qquad Ui=i,
\end{equation}
and
\begin{equation}\label{analytic:eq:perturbedcontraction}
 p_Fi=1,\qquad p_FD=0,\qquad Dh_F+h_FD=1-ip_F .
\end{equation}
All these operators are well defined, and applied to input support $S$ the
operators $U^{-1},p_F,h_F$ have output support in $S+E^*$.
\end{theorem}
\begin{proof}
The inverse of $U$ is supplied by \cref{analytic:lem:operator}. Since the
complex is concentrated in nonnegative degrees and $i(V)$ lies in degree
zero, $\delta i=0$. From $D^2=0$ we get $d\delta+\delta d+\delta^2=0$, and
\[
 Ud-DU=\delta hd-\delta-d\delta h-\delta^2h
 =\delta hd-\delta+\delta dh=\delta(hd+dh-1)=-\delta ip=0 ,
\]
so $Ud=DU$; and $Ui=i$ because $hi=0$, giving
\eqref{analytic:eq:chainconjugacy}. Since $h^2=0$ we have $Uh=h$ and hence
$UhU^{-1}=h_F$; conjugating $dh+hd=1-ip$ by $U$ gives the last identity in
\eqref{analytic:eq:perturbedcontraction}, and the other two follow from
$U^{-1}i=i$ and $U^{-1}D=dU^{-1}$. Every identity is an equality of
strongly summable coefficients, justified by
\cref{analytic:lem:neumann}; the support bound follows from the geometric
inverse for $U$ together with the fact that $p,h$ preserve support.
\end{proof}

The absence of a transferred differential on $V$ is not an extra
assumption: it follows from concentration in degree zero. In a complex with
homology in several degrees, positive perturbations can produce a
nontrivial differential on the homology space, and the theorem would have to
be modified.

\begin{lemma}[Ordinary contraction on one fixed neighbourhood]
\label{analytic:lem:ordinarykoszul}
Let $f_1,\ldots,f_n$ be ordinary holomorphic germs with an isolated common
zero at the origin and local algebra $A_0=R_n/(f)$ of dimension $\mu$, and
let $b_1,\ldots,b_\mu$ be polynomial representatives of a basis. Choose a
sufficiently small Stein polydisk $\Omega$ on which all $f_i$ are
holomorphic with no common zero except the origin, and form the
global-section Koszul complex
\begin{equation}\label{analytic:eq:koszul}
 C_k=\OO(\Omega)\otimes_\C\textstyle\bigwedge^k\C^n,\qquad
 d_f(g\,e_{i_1}\wedge\cdots\wedge e_{i_k})
 =\sum_{j=1}^k(-1)^{j-1}f_{i_j}g\,
 e_{i_1}\wedge\cdots\widehat{e_{i_j}}\cdots\wedge e_{i_k}.
\end{equation}
Then \eqref{analytic:eq:koszul} has no positive-degree homology and its
degree-zero homology is $A_0$; it therefore admits a contraction
\eqref{analytic:eq:contraction} with $V=\C^\mu$ and
$i(c_1,\ldots,c_\mu)=\sum_jc_jb_j$. Every component of its maps takes
holomorphic functions on $\Omega$ to holomorphic functions on that same
$\Omega$.
\end{lemma}
\begin{proof}
At the origin the $n$ germs generate an ideal of finite colength in a
regular local ring of dimension $n$, so they are a system of parameters and
a regular sequence and their Koszul complex is exact in positive degrees. At
any other point of $\Omega$ some $f_i$ is a unit in the stalk and the Koszul
complex is contractible. Thus the sheaf Koszul complex resolves a coherent
quotient sheaf supported at the origin with stalk $A_0$. Kernels in this
finite complex are coherent; Cartan's theorem~B on the Stein domain
$\Omega$ makes their higher cohomology vanish, and applying it to the
successive short exact sequences shows that global sections preserve the
required exactness, the global sections of the quotient sheaf being $A_0$.
These ordinary analytic and regular-local facts are classical
\cite{GunningRossi,GrauertRemmert,Stacks}. The vector-space splitting
described above now supplies the maps, with codomains the fixed spaces
\eqref{analytic:eq:koszul}.
\end{proof}

This use of Cartan~B is substantive, and it is exactly the step that has no
analogue over $\A_n$: choosing a separate germ-level division for each
coefficient and then intersecting all its neighbourhoods would not prove the
final assertion of the lemma.

\begin{theorem}[Support-explicit conservation of multiplicity over $\RR_n$]
\label{analytic:thm:conservationR}
Let $f,\mu,b_1,\ldots,b_\mu$ be as above, let $F_i=f_i+E_i$ with all
$E_i\in\RR_n$ of strictly positive Hahn support (zero errors allowed), and
put $E=\bigcup_i\supp E_i\subset\Gamma_{>0}$. Then:
\begin{enumerate}
\item the classes of $b_1,\ldots,b_\mu$ form a $K$-basis of
$A_F=\RR_n/(F_1,\ldots,F_n)$, so $\dim_KA_F=\mu$;
\item for $G\in\RR_n$ there are $K$-linear normal-form coordinates
$\NF_F(G)_j$ and division certificates $Q_i(G)\in\RR_n$ with
\begin{equation}\label{analytic:eq:normalformR}
 G=\sum_{j=1}^\mu b_j\,\NF_F(G)_j+\sum_{i=1}^nF_iQ_i(G),
\end{equation}
and on one fixed ordinary domain representing the inputs these outputs have
support in $\supp G+E^*$;
\item all higher homology of the Koszul complex of $F_1,\ldots,F_n$ over
$\RR_n$ vanishes;
\item the multiplication matrices $M_{z_\ell}$ in the prescribed basis have
entries supported in $E^*$, and their coefficients at exponent zero are the
ordinary multiplication matrices of $z_\ell$ on $A_0$.
\end{enumerate}
The normal-form coordinates are intrinsic once the basis is fixed; the
division certificates may depend on the ordinary contraction chosen in the
proof.
\end{theorem}
\begin{proof}
Take one small Stein polydisk $\Omega$ representing all errors, with the
ordinary zero set of $f$ equal to $\{0\}$ there, choose the contraction of
\cref{analytic:lem:ordinarykoszul}, and extend it coefficientwise to the
Hahn Koszul complex over $\HH_\Gamma(\Omega)$. The perturbation
$\delta=d_F-d_f$ is the Koszul differential formed from $E_1,\ldots,E_n$; it
has support in $E$ and each coefficient is a linear operator on the same
holomorphic function spaces. Koszul differentials square to zero, so
\cref{analytic:thm:conjugacy} applies. In degree zero set
\begin{equation}\label{analytic:eq:normalformoperator}
 \NF_F(G)=p(1+\delta h)^{-1}G,\qquad
 Q(G)=h(1+\delta h)^{-1}G\in C_1((t^\Gamma)).
\end{equation}
The contraction identity gives $G=i\,\NF_F(G)+d_FQ(G)$, which is
\eqref{analytic:eq:normalformR}, and $\NF_F\circ i=1$,
$\NF_F\circ d_F=0$; hence $\NF_F$ identifies the quotient with $K^\mu$ with
basis the classes of the $b_j$. The conjugacy of complexes gives the
positive-degree exactness. Every operation is performed on the one domain
$\Omega$, and \cref{analytic:thm:conjugacy} gives the support bound.

For an additional germ $G$, choose a smaller common Stein polydisk if
needed and repeat; restrictions of the quotient algebras send each $b_j$ to
the same $b_j$ and are isomorphisms because these are bases on both
domains, so the normal forms agree under restriction. Passing to the direct
limit proves the assertions over $\RR_n$ and their independence from the
chosen contraction, and positive-degree exactness likewise passes to the
limit. Finally the $j$th column of $M_{z_\ell}$ is $\NF_F(z_\ell b_j)$; its
input is ordinary, so its support is $\{0\}$ and the output support lies in
$E^*$, and at exponent zero the inverse in
\eqref{analytic:eq:normalformoperator} is the identity, so the column
reduces to $p(z_\ell b_j)$.
\end{proof}

\subsection{A coefficient recursion, not an iteration limit}
\label{analytic:subsec:recursion}
Formula \eqref{analytic:eq:normalformoperator} has an explicit reading. In
degree zero put $L=\delta h=\sum_{e\in E}t^eL_e$ and $W=(1+L)^{-1}G$; then
\begin{equation}\label{analytic:eq:Wrecursion}
 W_\nu=G_\nu-\!\!\sum_{\substack{e\in E,\ \eta\in\supp G+E^*\\ e+\eta=\nu}}
 \!\!L_eW_\eta ,
\end{equation}
where every $\eta$ on the right is smaller than $\nu$, and for fixed $\nu$
only finitely many input coefficients and finitely many words in the
perturbation coefficients are involved. Then $\NF_{F,\nu}=pW_\nu$ and
$Q_\nu=hW_\nu$. The analogous recursion over $\A_n$ is obtained from
\eqref{analytic:eq:normal-form} with $T$ in place of $L$.

This is an effective recursion \emph{relative to} available ordinary
division maps and an effective presentation of the support. Arbitrary
complex-linear splittings and arbitrary well-ordered sets need not come with
computable presentations, and the theorem does not assert a finite-time
enumeration algorithm for every set-theoretically possible Hahn input. The
finite-dependence statement at a specified coefficient, however, is
unconditional.

\begin{corollary}[Infinitesimal coordinate eigenvalues]
\label{analytic:cor:matrices}
Over either ring, the matrices $M_{z_1},\ldots,M_{z_n}$ commute, and for
each $\ell$
\begin{equation}\label{analytic:eq:characteristicreduction}
 \det(XI-M_{z_\ell})\equiv X^\mu \pmod{\m_K[X]} .
\end{equation}
Every coordinate eigenvalue of the deformed finite algebra is therefore
infinitesimal.
\end{corollary}
\begin{proof}
The matrices represent multiplication in a commutative algebra, so they
commute. In $A_0$ every coordinate lies in the nilpotent maximal ideal and
its multiplication matrix is nilpotent with characteristic polynomial
$X^\mu$; the reduction statement follows from the exponent-zero assertions
of \cref{analytic:thm:conservationR} over $\RR_n$ and of
\cref{analytic:subsec:matricesA} over $\A_n$. A root of a monic polynomial
whose lower coefficients are infinitesimal must be infinitesimal, by the
argument of \cref{analytic:cor:freehypersurfaceR}.
\end{proof}

\section{Residue duality, traces, and exact zero multiplicities}
\label{analytic:sec:trace}

\subsection{The contour-expansion functional over $\RR_n$}
\label{analytic:subsec:contour-residue}
Over $\RR_n$ one fixed ordinary residue cycle is available, and expanding
on it produces a functional whose construction is visibly independent of the
splitting used in \cref{analytic:thm:conservationR}. For an ordinary
isolated complete intersection $f$ in a small neighbourhood of the origin
choose a standard Grothendieck residue cycle
\[
 \mathcal C_f=\{z:|f_1(z)|=\epsilon_1,\ldots,|f_n(z)|=\epsilon_n\}
\]
in a suitable sufficiently small finite-map neighbourhood, the radii chosen
so that this is a compact regular real $n$-cycle oriented by
$d\arg f_1\wedge\cdots\wedge d\arg f_n>0$; equivalent cycles give the same
integral, and the normalization is
\eqref{analytic:eq:classical-residue}. On a neighbourhood of
$\mathcal C_f$ every $f_i$ is nonzero, so
\begin{equation}\label{analytic:eq:residueexpansion}
 \frac1{F_i}=\frac1{f_i}\sum_{k\ge0}\Bigl(-\frac{E_i}{f_i}\Bigr)^k
\end{equation}
is a legitimate common-support Hahn expansion there. For a uniformly
supported Hahn family of differential forms on a compact ordinary cycle,
define its $H$-integral coefficientwise; this is a coefficient functional,
and the cycle is not being treated as a fine-continuous parametrized object
in $\SC^n$. Define
\begin{equation}\label{analytic:eq:hahn-contour-residue}
 \Lambda_F(G)=\frac1{(2\pi i)^n}
 \HInt_{\mathcal C_f}
 \frac{G\,dz_1\wedge\cdots\wedge dz_n}{F_1\cdots F_n},
\end{equation}
the form being expanded by \eqref{analytic:eq:residueexpansion} before its
ordinary coefficients are integrated.

\begin{theorem}[Support-controlled residue duality over $\RR_n$]
\label{analytic:thm:residueduality}
For the data of \cref{analytic:thm:conservationR},
\eqref{analytic:eq:hahn-contour-residue} is well defined and independent of
the chosen sufficiently small ordinary residue cycle. It descends to a
$K$-linear functional $\Lambda_F:A_F\to K$ with
\begin{equation}\label{analytic:eq:residuesupport}
 \supp\Lambda_F(G)\subseteq\supp G+E^* ,
\end{equation}
and the pairing $(u,v)\mapsto\Lambda_F(uv)$ on $A_F\times A_F$ is perfect.
In the prescribed basis its Gram matrix $\mathcal G_F$ has support in $E^*$
and satisfies
\[
 \mathcal G_F=\mathcal G_0+(\text{positive-support terms}),\qquad
 \mathcal G_0=(\res_f(b_rb_s))_{r,s},\qquad \det\mathcal G_0\ne0 .
\]
All these assertions remain valid when roots collide and the algebra becomes
nonreduced.
\end{theorem}
\begin{proof}
Expanding all denominators gives coefficients supported in
$\supp G+E^*$ with finitely many contributions at each exponent, by
\cref{analytic:lem:neumann}; on the cycle each contribution is an ordinary
differential form, so the integral exists coefficientwise with the stated
support. The individual terms are ordinary Grothendieck residue expressions
with denominators $f_1^{k_1+1}\cdots f_n^{k_n+1}$, and their ordinary cycle
independence gives independence of
\eqref{analytic:eq:hahn-contour-residue}, including passage to smaller
common neighbourhoods.

For $G=F_jH$, cancel the factor $F_j$ before expanding: every resulting
coefficient is an ordinary form with no $f_j$ in its denominator, and its
integral vanishes, since after adjoining a factor $f_j$ to numerator and
denominator it is a residue for the complete intersection of the indicated
powers with numerator in its defining ideal. Thus $\Lambda_F$ annihilates
$(F_1,\ldots,F_n)$; scalar Hahn linearity follows from finite convolution at
each exponent.

For the Gram matrix, each $b_rb_s$ is ordinary, and at exponent zero the
errors in the denominator expansion disappear, leaving exactly the classical
matrix $\mathcal G_0$, whose determinant is nonzero by ordinary residue
nondegeneracy. Hence $\det\mathcal G_F$ has nonzero coefficient
$\det\mathcal G_0$ at exponent zero and $\mathcal G_F$ is invertible over
$K$; since the classes of the $b_r$ form a basis of $A_F$ by
\cref{analytic:thm:conservationR}, the pairing is perfect. The argument
never uses nonvanishing of a Jacobian at any root, which is the collision
assertion.
\end{proof}

\begin{proposition}[The two residue functionals agree where both are
defined]\label{analytic:prop:residues-agree}
Let $E_j\in\RR_n$, so that both $\Res_F$ of
\eqref{analytic:eq:hahn-residue} and $\Lambda_F$ of
\eqref{analytic:eq:hahn-contour-residue} are defined on $\RR_n$. Then
$\Res_F=\Lambda_F$.
\end{proposition}
\begin{proof}
Expand \eqref{analytic:eq:hahn-contour-residue} by
\eqref{analytic:eq:residueexpansion} and integrate coefficientwise. The
term indexed by $k\in\N^n$ is
$(-1)^{|k|}$ times the ordinary integral of
$G\,E_1^{k_1}\cdots E_n^{k_n}\,dz/(f_1^{k_1+1}\cdots f_n^{k_n+1})$ over
$\mathcal C_f$, that is $(-1)^{|k|}\res_{f^{k+1}}(GE^k)$ by
\eqref{analytic:eq:classical-residue}. Summing over $k$ is exactly
\eqref{analytic:eq:hahn-residue}, and the rearrangement is legitimate
because \cref{analytic:lem:neumann} makes the family strongly summable.
\end{proof}

Accordingly we write $\Res_F$ for both from now on, while remembering that
over $\A_n$ only the series definition exists and no ordinary cycle carries
all the coefficients.

\begin{corollary}[A support-controlled constant residue frame]
\label{analytic:cor:residueframe}
Let $C=\mathcal G_0^{-1}(\mathcal G_F-\mathcal G_0)$. The matrix series
\begin{equation}\label{analytic:eq:residueframe}
 S=(I+C)^{-1/2}=\sum_{k\ge0}\binom{-1/2}{k}C^k
\end{equation}
is strongly summable with support in $E^*$ and satisfies
$S^{\mathsf T}\mathcal G_FS=\mathcal G_0$. Thus the residue bilinear space
has a basis congruent to the prescribed ordinary basis in which its matrix
is exactly the ordinary residue matrix, even on the collision locus. The
statement holds over $\RR_n$ and, with $\mathcal G_F$ as in
\cref{analytic:thm:finite-deformA}, over $\A_n$.
\end{corollary}
\begin{proof}
The matrix $C$ has positive support, so \cref{analytic:lem:neumann} makes
every coefficient of \eqref{analytic:eq:residueframe} a finite sum, and
ordinary one-variable formal power-series identities in the single matrix
$C$ give $S^2(I+C)=I$. Symmetry of $\mathcal G_F$ and $\mathcal G_0$ gives
$C^{\mathsf T}\mathcal G_0=\mathcal G_0C$, hence
$S^{\mathsf T}\mathcal G_0=\mathcal G_0S$. Since
$\mathcal G_F=\mathcal G_0(I+C)$ and $S$ commutes with $C$, the isometry
follows.
\end{proof}

This identifies bilinear spaces, not algebras. It does not erase the
nonreduced algebraic structure of a collision; in particular a perfect
residue pairing should not be confused with the usually degenerate trace
pairing of a nonreduced finite algebra.

\subsection{Finite-parameter testing without a uniform radius}
\label{analytic:subsec:finite-dependency}
The passage from ordinary identities to a Hahn identity needs care over
$\A_n$, since an infinite perturbation cannot simply be assigned infinitely
many small complex parameters and placed on one common contour. The
following lemma is the substitute, and it is the exact place where
\cref{analytic:rem:initial-segments} is used.

\begin{lemma}[Finite dependency at a prescribed exponent]
\label{analytic:lem:finite-dependency}
Fix the leading system $f$, a positive perturbation-support set $E$, and a
well-ordered numerator support $S$. At any prescribed exponent $\gamma$,
every coefficient constructed from \eqref{analytic:eq:normal-form},
\eqref{analytic:eq:hahn-residue}, finite products and coefficientwise
differentiation depends on only finitely many perturbation coefficient
germs and finitely many numerator coefficient germs, and the contributing
words have a finite maximum length.
\end{lemma}
\begin{proof}
Each construction has an expansion whose contributions are indexed by a
source exponent in $S$ and a finite word in $E$. Finite products only
concatenate finitely many such words, and differentiation changes
coefficient germs but not their exponents. \Cref{analytic:lem:neumann}
gives finitely many words producing $\gamma$, including a finite maximum
length. Retaining exactly the source and perturbation coefficients occurring
in them determines the coefficient in question. No assertion is made that
there are only finitely many exponents below $\gamma$.
\end{proof}

\begin{theorem}[Trace--Jacobian identity]\label{analytic:thm:trace}
Let $B_F$ denote $\A_n/(F)$ under the hypotheses of
\cref{analytic:thm:finite-deformA}, or $\RR_n/(F)$ under those of
\cref{analytic:thm:conservationR}. In either case, for every $H$ in the
ambient ring,
\begin{equation}\label{analytic:eq:trace}
 \boxed{\qquad \Res_F(HJ_F)=\Tr_K(M_H),\qquad
 J_F=\det\Bigl(\frac{\partial F_i}{\partial z_j}\Bigr).\qquad}
\end{equation}
In particular $\Res_F(J_F)=\mu$.
\end{theorem}
\begin{proof}
Fix a Hahn exponent $\gamma$. By \cref{analytic:lem:finite-dependency} the
coefficients at $\gamma$ on both sides involve only finitely many ordinary
analytic coefficient germs. Temporarily replace the associated perturbation
monomials by independent ordinary complex parameters $u_1,\ldots,u_q$,
treating the finitely many numerator coefficients separately by linearity.
This produces a finite analytic deformation
$g_u=f+\sum_{\ell}u_\ell q_\ell$ of $f$ --- a finite family, whose germs do
have a common ordinary representative neighbourhood even when the original
Hahn family does not.

For that finite deformation, the ordinary cluster algebra is finite free
with the chosen basis after restricting the parameter neighbourhood, and
\eqref{analytic:eq:ordinary-trace} holds; the justification is spelled out
in \cref{analytic:app:classical-family}. The Taylor series in the
parameters of its cluster residue is the geometric denominator expansion
\eqref{analytic:eq:hahn-residue}, since on a sufficiently small ordinary
residue cycle the expansion converges uniformly. The formal parameter
expansion of its normal-coordinate map is
\eqref{analytic:eq:normal-form}, respectively
\eqref{analytic:eq:normalformoperator}: that Neumann expression gives a
decomposition in the parameter-adic completion, where the basis coordinates
are unique. This identifies the expansion even when the chosen linear
splitting $h$ was not norm-continuous.

Now substitute $u_\ell=t^{e_\ell}$ for the finitely many retained
perturbation exponents and extract the coefficient of $t^\gamma$. If several
parameter monomials carry the same Hahn exponent their contributions are
added, and there are finitely many by
\cref{analytic:lem:neumann}; omitted coefficient germs cannot contribute,
by their selection. This proves equality at $\gamma$, and $\gamma$ was
arbitrary. Taking $H=1$ gives the last assertion.
\end{proof}

The point of this proof is not to rebrand the ordinary trace identity as
new. It is to verify that the identity survives the absence of an
Archimedean valuation scale, and --- over $\A_n$ --- the absence of a common
ordinary radius.

\subsection{Actual points account for the whole length}
\label{analytic:subsec:zeros}
\begin{theorem}[Conservation of singular zeros]\label{analytic:thm:zeros}
Let $Z_F=\{a\in\m_K^n:F_1(a)=\cdots=F_n(a)=0\}$, over $\A_n$ under the
hypotheses of \cref{analytic:thm:finite-deformA} or over $\RR_n$ under
those of \cref{analytic:thm:conservationR}. Then $Z_F$ is finite and
nonempty, and there is an Artin-algebra decomposition
\begin{equation}\label{analytic:eq:artin}
 B_F\simeq\prod_{a\in Z_F}B_a,\qquad e_a=\dim_KB_a\ge1,\qquad
 \sum_{a\in Z_F}e_a=\mu ,
\end{equation}
with $e_a$ given by \eqref{analytic:eq:localmultiplicity}. The points of
$Z_F$ are exactly the joint characters of the commuting multiplication
matrices $M_{z_1},\ldots,M_{z_n}$, that is, their simultaneous eigenvalue
tuples. A point is simple, equivalently $e_a=1$, exactly when
$J_F(a)\ne0$.
\end{theorem}
\begin{proof}
$B_F$ is a nonzero finite-dimensional commutative $K$-algebra, hence
Artinian, and decomposes into its finitely many local factors. By
\cref{analytic:thm:weak-nullA} over $\A_n$, respectively
\cref{analytic:thm:weaknullR} over $\RR_n$, its maximal ideals are exactly
the common zero evaluations with residue field $K$; this gives
\eqref{analytic:eq:artin}. The factor at $a$ is the quotient of the local
ring at $\m_a$ by the ideal generated by the $F_i$; it has finite length, so
completion does not change it, and
\cref{analytic:thm:regularA}, respectively
\cref{analytic:prop:completionR}, identifies the completed quotient,
giving \eqref{analytic:eq:localmultiplicity}.

In the local factor at $a$, multiplication by $z_j$ is $a_j$ times the
identity plus a nilpotent operator, so simultaneous triangularization of the
commuting multiplication operators puts precisely these tuples on the
diagonal with the corresponding local lengths; conversely a common
eigenvector in the regular representation determines a character, and the
evaluation description identifies that character with a point of $Z_F$.

For the simple-point criterion, work in the formal local quotient: its
cotangent space is $K^n$ modulo the rows of the Jacobian matrix at $a$. If
that matrix is invertible the maximal ideal has zero cotangent space, hence
is zero by Nakayama's lemma and the local algebra is $K$; conversely a
length-one local algebra is $K$ with zero cotangent space, forcing the
Jacobian to have rank $n$.
\end{proof}

The theorem protects total \emph{length}. It neither requires nor implies
that the number of distinct points stays constant: a point of length $\mu$
can split into $\mu$ simple points or into intermediate nonreduced clusters.

\begin{proposition}[Four equivalent forms of the simple-root criterion]
\label{analytic:prop:jacobian}
Over either ring, the following are equivalent: every point of $Z_F$ has
local multiplicity one; $J_F(a)\ne0$ for all $a\in Z_F$; $J_F$ is a unit in
$B_F$; the multiplication matrix $M_{J_F}$ has nonzero determinant. When
they hold, $Z_F$ has exactly $\mu$ distinct points.
\end{proposition}
\begin{proof}
The first equivalence is the last paragraph of
\cref{analytic:thm:zeros}. An element of an Artinian algebra is a unit
exactly when it lies in no maximal ideal, which gives the second;
multiplication by an element of a finite-dimensional algebra is invertible
exactly when that element is a unit, which gives the third.
\end{proof}

\begin{corollary}[Weighted evaluation and simple residues]
\label{analytic:cor:weighted}
For every $H$ in the ambient ring,
\begin{equation}\label{analytic:eq:weighted}
 \Res_F(HJ_F)=\sum_{a\in Z_F}e_a\,H(a),
\end{equation}
and if all points are simple then
\begin{equation}\label{analytic:eq:simple-residue}
 \Res_F(H)=\sum_{a\in Z_F}\frac{H(a)}{J_F(a)} .
\end{equation}
\end{corollary}
\begin{proof}
On $B_a$, multiplication by $H$ is the scalar $H(a)$ plus a nilpotent
operator, so its trace is $e_aH(a)$; summing and applying
\cref{analytic:thm:trace} gives \eqref{analytic:eq:weighted}. If every point
is simple then $B_F\simeq K^{Z_F}$; letting $\epsilon_a$ be the
corresponding idempotent, \eqref{analytic:eq:trace} applied to
$\epsilon_a$ gives $\Res_F(\epsilon_a)J_F(a)=1$, and expanding the class of
$H$ as $\sum_aH(a)\epsilon_a$ gives
\eqref{analytic:eq:simple-residue}. Alternatively, apply
\eqref{analytic:eq:trace} to $H/J_F$, which is legitimate by
\cref{analytic:prop:jacobian}.
\end{proof}

Individual summands in \eqref{analytic:eq:simple-residue} may be infinite
surcomplex numbers; their finite sum is nevertheless the strongly defined
residue \eqref{analytic:eq:hahn-residue}, and the support of the total can
be far smaller than the supports of the individual local contributions.
\Cref{analytic:sec:example} makes this cancellation explicit twice.

\begin{remark}[What is and is not localized]
\label{analytic:rem:localization}
At a multiple point \eqref{analytic:eq:simple-residue} is not a valid
formula with a zero denominator. The always-valid statement is the
functional on the finite local algebra: the Artinian product decomposition
restricts $\Res_F$ to a perfect functional on each local factor, their sum
is the total functional, and the trace--Jacobian identity still holds. This
gives an algebraic local residue interpretation without inventing small
fine-topological integration cycles around individual nonstandard points.
\end{remark}

\section{Workspace invariance and the analytic--formal comparison}
\label{analytic:sec:extension}

\subsection{No additional roots at unrepresented surreal scales}
\label{analytic:subsec:workspace-invariance}
Let $\Gamma\subseteq\Delta\subset\No$ be nonzero divisible set-sized
ordered subgroups, and use subscripts to distinguish
$K_\Gamma\subseteq K_\Delta$, $\A_{n,\Gamma}\subseteq\A_{n,\Delta}$ and
$\RR_n(\Gamma)\subseteq\RR_n(\Delta)$. Both source manuscripts prove a
result of this shape, over their own rings and by different arguments; we
state both and compare the proofs rather than duplicating one of them.

\begin{theorem}[Workspace invariance over $\A_n$]
\label{analytic:thm:workspace}
For a system $F$ as in \cref{analytic:thm:finite-deformA} with coefficients
in $\A_{n,\Gamma}$ there is a natural isomorphism
\begin{equation}\label{analytic:eq:base-change}
 \frac{\A_{n,\Delta}}{(F)}\ \simeq\
 K_\Delta\otimes_{K_\Gamma}\frac{\A_{n,\Gamma}}{(F)} .
\end{equation}
The infinitesimal zero set is already contained in $K_\Gamma^n$ and its
local multiplicities do not change over $K_\Delta$.
\end{theorem}
\begin{proof}
Use the same ordinary basis $b_1,\ldots,b_\mu$ and the same splitting
$p,h$ in both workspaces. The proof of
\cref{analytic:thm:finite-deformA} applies in each and gives this basis on
each side. The multiplication constants computed by
\eqref{analytic:eq:normal-form} for the products $b_rb_s$ already lie in
$K_\Gamma$ and are unchanged in $K_\Delta$, so the canonical map in
\eqref{analytic:eq:base-change} identifies the bases and their products.

A point over $K_\Delta$ defines a character of this finite algebra; every
coordinate of the point satisfies the characteristic polynomial of the
corresponding coordinate multiplication matrix over $K_\Gamma$, and since
$K_\Gamma$ is algebraically closed all roots of that polynomial lie in
$K_\Gamma$. Hence all coordinates already lie in the smaller field.
Equivalently, base change of each finite local factor retains its residue
field and its dimension, the nilpotent maximal ideal remaining nilpotent,
which is preservation of points and multiplicities.
\end{proof}

\begin{theorem}[Base change and full surcomplex exhaustiveness over $\RR_n$]
\label{analytic:thm:basechange}
For the data of \cref{analytic:thm:conservationR} in $\RR_n(\Gamma)$, with
$K'=K_\Delta$, the natural map
\begin{equation}\label{analytic:eq:basechangeR}
 A_F(\Gamma)\otimes_KK'\ \xrightarrow{\ \sim\ }\ A_F(\Delta)
\end{equation}
is an isomorphism. In the prescribed polynomial basis the multiplication
matrices on both sides are \emph{identical}. The zero set in $\m_{K'}^n$ is
exactly the original zero set in $\m_K^n$ with unchanged local
multiplicities. Consequently, evaluating the same coherent system on the
entire surcomplex monad introduces no additional roots, and the count
\eqref{analytic:eq:artin} is valid there as well.
\end{theorem}
\begin{proof}
The construction \eqref{analytic:eq:normalformoperator} depends only on the
ordinary contraction and on the existing perturbation coefficients;
enlarging the exponent group changes no coefficient calculation. By
\cref{analytic:thm:conservationR} the same $b_1,\ldots,b_\mu$ form a basis
on each side of \eqref{analytic:eq:basechangeR}, and the natural map sends
one basis to the other, hence is an algebra isomorphism. The characteristic
polynomial of each multiplication matrix has coefficients in $K$ and all its
roots already lie in $K$ by algebraic closedness, so any coordinate of a
new common zero in $K'$ would be one of these roots; evaluation then agrees
with the original evaluation and no new point appears. Base change of the
Artinian local factors preserves their dimensions and hence their lengths.
Finally, the finitely many coordinates of any proposed surcomplex point have
set-sized supports; enlarging $\Gamma$ by those supports and taking the
divisible subgroup they generate gives a set-sized $\Delta$ to which the
preceding argument applies, by \cref{analytic:lem:enlargement}.
\end{proof}

\begin{corollary}[Full-surcomplex conservation]
\label{analytic:cor:full-sc}
Interpret the equations of \cref{analytic:thm:workspace} at arbitrary
surcomplex tuples infinitesimal relative to $\C$. Every such common zero
belongs to the original field $K_\Gamma^n$. Consequently the exact count
$\sum_ae_a=\mu$ already accounts for \emph{all} full-surcomplex
infinitesimal zeros.
\end{corollary}
\begin{proof}
Any proposed tuple lies in a set-sized divisible workspace $K_\Delta$
containing $K_\Gamma$, by \cref{analytic:lem:enlargement}, and evaluation in
it is compatible with the smaller workspace by the strong-sum definitions.
Apply \cref{analytic:thm:workspace}.
\end{proof}

\begin{remark}[The two proofs, compared]
\label{analytic:rem:two-workspace-proofs}
\Cref{analytic:thm:workspace} runs the $\A_n$ deformation argument twice,
once in each workspace, and identifies the two resulting bases together with
their multiplication constants; the base-change isomorphism is then a
consequence of having the same basis on both sides.
\Cref{analytic:thm:basechange} instead observes that the single
construction \eqref{analytic:eq:normalformoperator} is literally unchanged
by enlarging $\Gamma$, because it reads only the ordinary contraction and
the existing perturbation coefficients, and therefore yields the
\emph{identical} matrices rather than matched ones. The second formulation
is the sharper statement about matrices; the first is available over the
strictly larger ring. Both then close the argument the same way, by
algebraic closedness of the smaller field. Neither is a corollary of the
other, and neither may be transferred between rings.
\end{remark}

This rules out a specific possible defect of a set-sized construction: an
isolated cluster cannot acquire extra roots merely because a new surreal
scale is allowed. The conclusion uses both finite dimensionality and
algebraic closedness, and it is not a statement that arbitrary
infinite-dimensional equations have workspace-independent solution sets.
One further caution: the supports of individual root coordinates need not
lie in $E^*$. Fractional scales such as $t^{1/2}$ may appear in individual
roots even when every matrix coefficient uses only integral powers;
divisibility of $\Gamma$ accommodates these scales, and
\cref{analytic:subsec:radius-free-example} exhibits them.

\subsection{Allowing divergent ordinary coefficient series}
\label{analytic:subsec:formal}
There is also a useful comparison with the formal-coefficient ring
$\FF_n=\C[[z_1,\ldots,z_n]]((t^\Gamma))$, whose coefficients are arbitrary
formal complex power series. Strong evaluation still makes sense at every
infinitesimal tuple, but coefficientwise ordinary analyticity has been
dropped. The following theorem separates that change precisely from
completion at a point. It is a statement about $\A_n$ and $\FF_n$ only; no
analogue for $\RR_n$ is claimed.

\begin{theorem}[Analytic--formal comparison]
\label{analytic:thm:formal-comparison}
The inclusion $\A_n\hookrightarrow\FF_n$ is faithfully flat. Both rings have
the same evaluation maximal ideals, the same residue fields, and the same
completed local rings at corresponding points. If $I\subseteq\A_n$ has
finite-dimensional quotient over $K$, the natural map
\begin{equation}\label{analytic:eq:formal-finite}
 \A_n/I\longrightarrow\FF_n/I\FF_n
\end{equation}
is an isomorphism. In particular, allowing divergent ordinary coefficient
germs does not change a finite zero scheme defined over $\A_n$.
\end{theorem}
\begin{proof}
The arguments of
\crefrange{analytic:sec:germs}{analytic:sec:structure-A} work with $\C[[z]]$
in place of $\C\{z\}$, using formal Weierstrass division; translation,
finite jets and faithful evaluation use only formal Taylor identities and
the support lemma. Consequently $\FF_n$ is Noetherian, its maximal ideals
are the same coordinate evaluation ideals, and its local completion at $a$
is $K[[X]]$.

Let $A=(\A_n)_{\m_a}$ and $B=(\FF_n)_{\m_a}$. The induced map of
completions is the identity on $K[[X]]$, and completion of a Noetherian
local ring is faithfully flat \cite[tag 00MC]{Stacks}. It follows that
$A\to B$ is flat: tensoring a flatness obstruction
$\operatorname{Tor}_1^A(M,B)$ with the faithfully flat $B$-module
$\widehat B$ turns it into
$\operatorname{Tor}_1^A(M,\widehat B)=\operatorname{Tor}_1^A(M,\widehat A)=0$,
and faithfulness of $B\to\widehat B$ forces the obstruction to vanish.
Flatness may be checked at maximal ideals of the target, so
$\A_n\to\FF_n$ is flat, and every maximal ideal of the source has its
corresponding evaluation ideal in the target, making the map faithfully
flat.

Now assume $\dim_K\A_n/I<\infty$. Its Artin decomposition has finitely many
supporting points $a_1,\ldots,a_r$, and for some $N$ the product
$J=\prod_{\ell=1}^r\m_{a_\ell}^N$ is contained in $I$. The ideals for
distinct points are comaximal, so by the Chinese remainder theorem and the
identical finite-jet quotients of the two rings the map
$\A_n/J\to\FF_n/J\FF_n$ is an isomorphism; quotienting by the image of
$I/J$ gives \eqref{analytic:eq:formal-finite}.
\end{proof}

By \cref{analytic:rem:formal-strict} the inclusion is nevertheless strict,
so the comparison has content beyond a change of notation. Together with
\cref{analytic:ex:no-radius} it completes the picture promised in
\cref{analytic:subsec:four-rings}: the passage from common-domain to
radius-free coefficients is a genuine enlargement that no common-domain
theorem covers, whereas the further passage to formal coefficients,
although also a genuine enlargement of the ring, changes nothing about
finite-colength quotients. The comparison also provides a precise bridge to
formal or full separated Hahn analytic structures without asserting that the
two coefficient rings are equal.

\section{Quantitative Hensel--Rouch\'e stability over $\A_n$}
\label{analytic:sec:stability}

Conservation of total length does not by itself label roots after a
perturbation: a singular point may split, and several simple points may
collide. This section obtains a root-by-root statement over $\A_n$ with an
explicit threshold, formulated through the valuation loss in the inverse
Jacobian, which can be substantially better than a bound using only its
determinant.

\subsection{A local estimate at a simple root}
Assume $F\in\A_n^n$ with $\vH(F_i)\ge0$ and $F(a)=0$ for $a\in\m_K^n$. If
$L=DF(a)$ is invertible, define
\begin{equation}\label{analytic:eq:kappa}
 \kappa_a=\max\Bigl(0,-\min_{i,j}v\bigl((L^{-1})_{ij}\bigr)\Bigr),
\end{equation}
so every entry of $L^{-1}$ has valuation at least $-\kappa_a$, while every
entry of $L$ has nonnegative valuation. The adjugate formula yields the
convenient but \emph{sometimes strictly weaker} bound
\begin{equation}\label{analytic:eq:det-bound}
 0\le\kappa_a\le v(\det L).
\end{equation}
The inequality lives in $\Gamma$, an arbitrary ordered group; it is not an
estimate obtained by replacing valuations with real-valued magnitudes.

\begin{lemma}[Scaled Taylor polynomials at each Hahn exponent]
\label{analytic:lem:scaled-polynomial}
For $a\in\m_K^n$, $s\in\Gamma_{>0}$ and $F\in\A_n$, the expansion
$F(a+t^sy)$ has polynomial complex coefficient functions in $y$ at every
fixed Hahn exponent. Consequently it may be translated by any fixed complex
vector in the $y$ variables.
\end{lemma}
\begin{proof}
First translate by $a$ using \cref{analytic:prop:translation}, then expand
in monomials in $t^sy$. At a fixed exponent $\gamma$ the possible Taylor
degrees and source exponents satisfy $\lambda+rs=\gamma$, and
\cref{analytic:lem:neumann} gives finitely many such decompositions, hence
finitely many degrees; each degree has finitely many multivariate
monomials. The resulting coefficient is therefore a polynomial, so a fixed
complex translation is legitimate coefficientwise and preserves Hahn
support.
\end{proof}

\begin{theorem}[One-root Hensel--Rouch\'e bound over $\A_n$]
\label{analytic:thm:local-stability}
Under the hypotheses above let $G=F+D$ with $D\in\A_n^n$ nonzero and
$\delta=\min_i\vH(D_i)>2\kappa_a$. Then $G$ has a unique zero $b$ in the
ball
\begin{equation}\label{analytic:eq:stability-ball}
 \mathcal B_a=\{b\in\m_K^n:\vmin(b-a)>\kappa_a\},
\end{equation}
it satisfies
\begin{equation}\label{analytic:eq:displacement}
 \vmin(b-a)\ \ge\ \delta-\kappa_a ,
\end{equation}
and it is simple.
\end{theorem}
\begin{proof}
Put $\kappa=\kappa_a$ and $s=\delta-\kappa>\kappa$, and introduce the
normalized system $H(y)=t^{-s}L^{-1}G(a+t^sy)$. Its linear contribution from
$F$ is exactly $y$. Since all Taylor coefficients of $F$ at $a$ have
valuation at least zero, the terms of degree at least two in the normalized
contribution have coefficients of valuation at least
$2s-s-\kappa=s-\kappa=\delta-2\kappa>0$. The constant vector
$c=t^{-s}L^{-1}D(a)$ has valuation at least zero, and every nonconstant
contribution of $D$ has coefficient valuation at least $\delta-\kappa=s>0$.
Thus
\begin{equation}\label{analytic:eq:normalized-stability}
 H(y)=y+c_0+P(y),\qquad c_0=\red(c)\in\C^n,\qquad\vH(P_i)>0 .
\end{equation}
By \cref{analytic:lem:scaled-polynomial} the coefficient functions here are
polynomials, so the substitution $y=-c_0+u$ is valid even though the
original coefficients were only germs; the resulting system is
$u+\widetilde P(u)$ with strictly positive perturbation support. Applying
\cref{analytic:thm:finite-deformA,analytic:thm:zeros} to the ordinary
leading system $(u_1,\ldots,u_n)$, of length one, gives exactly one
infinitesimal zero $u$. Hence $b=a+t^s(-c_0+u)$ is a zero of $G$, lies in
$\m_K^n$ and satisfies \eqref{analytic:eq:displacement}; the length-one
assertion makes it simple, since the affine change of variables and the
invertible normalizing matrix preserve nonvanishing of the Jacobian.

It remains to prove uniqueness in the whole ball
\eqref{analytic:eq:stability-ball}, not merely in the residue class used to
construct the zero. Suppose $a+h$ and $a+k$ are zeros with
$\vmin(h),\vmin(k)>\kappa$. Formal Taylor subtraction, factoring each
difference of monomials, gives $G(a+h)-G(a+k)=(L+C)(h-k)$ where every entry
of $C$ has valuation at least $\min\{\vmin(h),\vmin(k),\delta\}>\kappa$:
terms from $F$ beyond its linear term retain at least one factor from $h$ or
$k$ after subtraction, and terms from $D$ retain valuation at least
$\delta$. All sums are strongly defined by
\cref{analytic:lem:neumann}. Thus $L^{-1}C$ has entries of strictly positive
valuation, so $I+L^{-1}C$ is invertible by strong matrix geometric
inversion, and the difference equation forces $h=k$.
\end{proof}

The construction proves existence inside $K$ itself. It does not invoke
sequential completeness or a rank-one Banach contraction theorem, and all
normalizations use actual monomials $t^s$ with $s\in\Gamma$.

\subsection{Correspondence for an entire finite cluster}
\begin{corollary}[Separation of simple zeros]
\label{analytic:cor:separation}
If $a\ne a'$ are two simple infinitesimal zeros of the same
nonnegative-support system $F\in\A_n^n$, then
\begin{equation}\label{analytic:eq:separation}
 \vmin(a-a')\le\min(\kappa_a,\kappa_{a'}),
\end{equation}
and the balls $\mathcal B_a$ around the simple zeros are pairwise disjoint.
\end{corollary}
\begin{proof}
The uniqueness argument of \cref{analytic:thm:local-stability} with $D=0$
says $a$ is the only zero of $F$ with $\vmin(b-a)>\kappa_a$; applying it at
both points gives \eqref{analytic:eq:separation}. If a point lay in both
balls the ultrametric inequality would give
$\vmin(a-a')>\min(\kappa_a,\kappa_{a'})$, a contradiction.
\end{proof}

\begin{theorem}[Finite-cluster correspondence]
\label{analytic:thm:cluster-stability}
Let $F$ satisfy \cref{analytic:thm:finite-deformA} with all $\mu$ zeros
simple, and let $G=F+D$ with
\begin{equation}\label{analytic:eq:global-threshold}
 \delta=\min_i\vH(D_i)>2\max_{a\in Z_F}\kappa_a .
\end{equation}
Then $G$ also has $\mu$ simple infinitesimal zeros, and there is a unique
bijection $a\mapsto b_a$ from $Z_F$ to $Z_G$ with $\vmin(b_a-a)>\kappa_a$
and $\vmin(b_a-a)\ge\delta-\kappa_a$. The same assertion accounts for all
infinitesimal surcomplex zeros.
\end{theorem}
\begin{proof}
The case $D=0$ is immediate. Otherwise each ball contains a unique simple
zero of $G$ by \cref{analytic:thm:local-stability}, and the balls are
disjoint by \cref{analytic:cor:separation}, so $\mu$ distinct simple zeros
have been found. Since $G$ is still a positive-support perturbation of the
same leading system $f$, \cref{analytic:thm:zeros} gives total length $\mu$,
so there are no others. The full-surcomplex assertion follows from
\cref{analytic:cor:full-sc}, applied in a workspace containing both
systems.
\end{proof}

\begin{example}[Sharpness of the strict threshold]
\label{analytic:ex:sharpness}
Fix $\eta>0$ and set $F(z)=z^2-t^{2\eta}$, $D(z)=t^{2\eta}$, $G(z)=z^2$.
The two roots of $F$ are $a=\pm t^\eta$ and $\kappa_a=\eta$. Here
$\delta=2\eta=2\kappa_a$, yet the two simple roots coalesce into a double
root, and the displacement to zero has valuation exactly $\eta$, not greater
than $\kappa_a$. Thus replacing the strict inequality in
\eqref{analytic:eq:global-threshold} by equality is false, and any universal
factor smaller than two is false as well. This is the familiar sharp scalar
Hensel obstruction, situated here in the analytic setting.
\end{example}

For higher-rank groups the threshold is an inequality in $\Gamma$, not an
estimate after replacing valuation by a real-valued magnitude. With
$\Gamma=\Q+\Q\omega$, a perturbation of valuation $N$ for arbitrarily large
natural $N$ still need not be smaller than a separation scale governed by
$\omega$.

\begin{remark}[Two certificates of stability, and which is sharper for
which purpose]\label{analytic:rem:two-certificates}
The material of this section straddles two reports and is stated in full in
the companion report on finite deformations \cite{FiniteDeformations},
which carries the complementary certificate as well. We record the
comparison here because \cref{analytic:thm:local-stability} is proved over
$\A_n$ and belongs to this report, and because averaging the two available
certificates into a single ``discriminant or Jacobian threshold'' would lose
real content.
\begin{enumerate}
\item The inverse-Jacobian certificate above is the one to use when
individual roots are to be \emph{matched}: it supplies the bijection
$a\mapsto b_a$ of \cref{analytic:thm:cluster-stability} together with the
displacement bound \eqref{analytic:eq:displacement}. A companion manuscript
proves in addition that both valuation bounds are attained, including the
second-order bound $v(h+A^{-1}\Delta(a))\ge2\sigma-3\kappa$; that
attainment statement is proved there and is not re-proved here.
\item A second certificate in that cluster \emph{declines root matching
altogether}, on the explicit ground that a labelling of individual roots
need not be Lipschitz, that roots of a multiple leading system generally
involve fractional exponents, and that inverse Jacobians can have negative
valuation. It certifies stability instead through a Discriminant--Jacobian
identity and a discriminant threshold, offering ``one certified geometric
consequence without choosing root branches.'' That certificate is the right
one when only the statement that all zeros remain simple is needed and no
root labelling is available.
\item Within the present section, \eqref{analytic:eq:det-bound} shows the
determinant-only bound is provably lossy:
\cref{analytic:subsec:four-roots} exhibits a system where
$\kappa_a=\max(\alpha,\beta)$ while $v(\det L)=\alpha+\beta$, so the
threshold $\delta>2\max(\alpha,\beta)$ is strictly better than
$\delta>2(\alpha+\beta)$.
\end{enumerate}
These are complementary certificates of different things, not competing
estimates of one thing.
\end{remark}

\section{Worked examples}\label{analytic:sec:example}

The first two examples are systems over $\A_2$; the second of them lies
outside $\RR_2$ and is the point of the radius-free ring. The third and
fourth are systems over $\RR_2$ and exhibit a collision and a
nonpolynomial, non-grid support.

\subsection{Four coupled roots and a perfect residue matrix, over $\A_2$}
\label{analytic:subsec:four-roots}
Let $A,B\in\m_K$ be nonzero and consider
\begin{equation}\label{analytic:eq:four-system}
 F_1=x^2-Ay,\qquad F_2=y^2-Bx .
\end{equation}
The ordinary leading ideal is $(x^2,y^2)$ with basis $1,x,y,xy$ and length
four, and \cref{analytic:thm:finite-deformA} gives the same basis for the
perturbed quotient without eliminating variables. Put $q=AB$; in the
quotient
\[
 x^2=Ay,\quad y^2=Bx,\quad x^2y=qx,\quad xy^2=qy,\quad (xy)^2=qxy .
\]
With columns indexed by multiplication of the ordered basis $(1,x,y,xy)$,
\begin{equation}\label{analytic:eq:example-matrices}
 M_x=\begin{pmatrix}0&0&0&0\\1&0&0&q\\0&A&0&0\\0&0&1&0\end{pmatrix},
 \qquad
 M_y=\begin{pmatrix}0&0&0&0\\0&0&B&0\\1&0&0&q\\0&1&0&0\end{pmatrix}.
\end{equation}
They commute and satisfy the two defining relations, and
\begin{equation}\label{analytic:eq:example-char}
 \chi_x(\lambda)=\lambda(\lambda^3-A^2B),\qquad
 \chi_y(\lambda)=\lambda(\lambda^3-AB^2).
\end{equation}
The Hahn residue expansion is explicit. From
\eqref{analytic:eq:hahn-residue} and
\eqref{analytic:eq:coordinate-residue},
\begin{equation}\label{analytic:eq:example-coeff}
 \Res_F(H)=\sum_{k,\ell\ge0}A^kB^\ell\,[x^{2k+1-\ell}y^{2\ell+1-k}]H ,
\end{equation}
a coefficient with a negative index being zero; the two signs from the
negative perturbations cancel the factor $(-1)^{k+\ell}$. For a monomial
$x^py^r$ the only possible indices are $k=(2p+r-3)/3$ and
$\ell=(p+2r-3)/3$, so its residue is $A^kB^\ell$ when these are nonnegative
integers and zero otherwise. In particular
$\Res_F(1)=\Res_F(x)=\Res_F(y)=0$ and $\Res_F(xy)=1$, whence
\begin{equation}\label{analytic:eq:example-gram}
 \mathcal G_F=\begin{pmatrix}
 0&0&0&1\\0&0&1&0\\0&1&0&0\\1&0&0&q\end{pmatrix},
 \qquad\det\mathcal G_F=1 .
\end{equation}
Its determinant remains a unit even as the points collide under
specialization of the parameters. This is why the residue pairing, rather
than a list of distinct roots, is the stable invariant.

Choose $r\in K$ with $r^3=A^2B$ and let $\zeta\in\C$ be a primitive cube
root of unity. The four roots are
\begin{equation}\label{analytic:eq:four-roots}
 (0,0),\qquad\Bigl(r\zeta^j,\frac{r^2}{A}\zeta^{2j}\Bigr),\quad j=0,1,2 ,
\end{equation}
since a nonzero solution has $y=x^2/A$ and hence $x^3=A^2B$, and conversely.
Writing $\alpha=v(A)>0$ and $\beta=v(B)>0$, the nonzero coordinate
valuations are
\begin{equation}\label{analytic:eq:four-valuations}
 v(x)=\frac{2\alpha+\beta}{3},\qquad v(y)=\frac{\alpha+2\beta}{3},
\end{equation}
all infinitesimal, with no assumption that $\alpha$ and $\beta$ are
Archimedean-equivalent. The Jacobian is $J_F=4xy-q$, equal to $-q$ at the
origin and to $3q$ at every nonzero root, so all four points are simple.
Their residue weights for $H=1$ are $-1/q,\ 1/(3q),\ 1/(3q),\ 1/(3q)$ and
sum to zero, each being infinite relative to $\C$; for $H=xy$ the sum is
$1$, and $\Res_F(J_F)=4$ while $\Tr(M_{xy})=3q$, in agreement with
\eqref{analytic:eq:example-coeff} and
\cref{analytic:thm:trace}. Taking $\Gamma=\Q+\Q\omega$, $A=t$ and
$B=t^\omega$, the nonzero coordinates have valuations $(2+\omega)/3$ and
$(1+2\omega)/3$ and every nonzero residue weight has valuation
$-(1+\omega)$; replacing this value group by a single real-valued small
parameter would discard the fact that $\omega$ dominates every positive
integer.

The inverse-Jacobian loss is computable exactly here. At zero
\[
 DF(0)^{-1}=\begin{pmatrix}0&-B^{-1}\\-A^{-1}&0\end{pmatrix},
 \qquad\text{and at a nonzero root}\qquad
 DF(a)^{-1}=\frac1{3AB}\begin{pmatrix}2y&A\\B&2x\end{pmatrix}.
\]
With \eqref{analytic:eq:four-valuations} this gives, at all four roots,
\begin{equation}\label{analytic:eq:example-kappa}
 \kappa_a=\max(\alpha,\beta),\qquad v(J_F(a))=\alpha+\beta .
\end{equation}
Thus a further perturbation with $\delta>2\max(\alpha,\beta)$ preserves all
four roots individually with the displacement bound of
\cref{analytic:thm:cluster-stability}, and the inverse-matrix threshold is
strictly better than the determinant-only sufficient condition
$\delta>2(\alpha+\beta)$, as promised in
\cref{analytic:rem:two-certificates}.

\subsection{A genuinely radius-free, higher-rank deformation over $\A_2$}
\label{analytic:subsec:radius-free-example}
Take $\Gamma=\Q+\Q\omega$ and consider the analytic-germ system
\begin{equation}\label{analytic:eq:radius-free-system}
 F_1=x^2+\sum_{r\ge1}\frac{t^r}{1-rx-r^2y},\qquad
 F_2=y^2+t^\omega x .
\end{equation}
The family of coefficient germs in $F_1$ has \emph{no} common ordinary
neighbourhood of the origin: restricted to $y=0$ it has poles at $x=1/r$.
Hence \eqref{analytic:eq:radius-free-system} is a system over $\A_2$ that
lies in no $\HH_\Gamma(U)$, and
\cref{analytic:thm:conservationR} does not apply to it; this is precisely
the gap that \cref{analytic:thm:finite-deformA} fills. Its support is the
well-ordered set of positive integers and the combined perturbation support
lies in $\N_{>0}\cup\{\omega\}$, so
\cref{analytic:thm:finite-deformA} applies and gives total length four on
the full surcomplex infinitesimal neighbourhood, by
\cref{analytic:cor:full-sc}.

More can be read off without solving the infinite equations. Make the
diagonal rescaling
\begin{equation}\label{analytic:eq:radius-free-scaling}
 x=t^{1/2}u,\qquad y=t^{\omega/2+1/4}v ,
\end{equation}
divide the first equation by $t$ and the second by $t^{\omega+1/2}$. The
resulting leading system is
\begin{equation}\label{analytic:eq:radius-free-leading}
 u^2+1=0,\qquad v^2+u=0 ,
\end{equation}
and every remaining coefficient has strictly positive Hahn valuation. After
this rescaling each Hahn coefficient is a polynomial in $(u,v)$, by the
coordinatewise version of
\cref{analytic:lem:scaled-polynomial}, so we may translate to each of the
four ordinary solutions $(u_0,v_0)$ of
\eqref{analytic:eq:radius-free-leading}. At each one
\[
 u_0=\pm i,\qquad v_0^2=-u_0,\qquad
 \det D(u^2+1,v^2+u)=4u_0v_0\ne0 ,
\]
so the length-one case of \cref{analytic:thm:finite-deformA} supplies one
infinitesimal lift in each translated chart; the charts are disjoint by
their different complex reductions, and since the unscaled system has total
length four these are all the zeros and all are simple. Their leading
coordinates are
\begin{equation}\label{analytic:eq:radius-free-roots}
 x=u_0t^{1/2}+\cdots,\qquad y=v_0t^{\omega/2+1/4}+\cdots ,
\end{equation}
with larger-valuation corrections, and at such a root
\begin{equation}\label{analytic:eq:radius-free-J}
 J_F=4u_0v_0\,t^{\omega/2+3/4}+\cdots :
\end{equation}
the derivative of the positive perturbation in $F_1$ has valuation at least
$1$, whereas $v(x)=1/2$ and $v(y)=\omega/2+1/4$, so $4xy$ dominates the
remaining terms of the determinant. Note that the root supports involve
$t^{1/2}$ and $t^{\omega/2+1/4}$, which lie in no monoid generated by the
perturbation support; this is the caution recorded after
\cref{analytic:rem:two-workspace-proofs}.

Although every individual simple residue $1/J_F(a)$ is infinite, their sum
is not merely finite --- it is infinitesimal:
\begin{equation}\label{analytic:eq:radius-free-residue}
 \boxed{\qquad\Res_F(1)=-4t+R,\qquad v(R)>1.\qquad}
\end{equation}
For this computation use $f=(x^2,y^2)$ in
\eqref{analytic:eq:hahn-residue}. The exponent-zero contribution vanishes.
A contribution at exponent $1$ can come only from $k=1$, $\ell=0$ and the
coefficient $1/(1-x-y)$ of $t$ in the first perturbation, and it equals
\[
 -[x^3y]\frac1{1-x-y}=-4 .
\]
Every other nonzero contribution has valuation strictly greater than $1$;
terms involving the second perturbation have valuation at least $\omega$.
Equation \eqref{analytic:eq:simple-residue} identifies
\eqref{analytic:eq:radius-free-residue} with the sum over the four infinite
local weights. This single example simultaneously uses shrinking ordinary
radii, several-variable coupling, different valuation ranks, and
cancellation beyond the leading pole scale.

\subsection{A length-five collision with three scale regimes, over $\RR_2$}
\label{analytic:subsec:length-five}
Let $s,u\in\m_K$ and consider
\begin{equation}\label{analytic:eq:exampleF}
 F_1=x^2-y^3-s,\qquad F_2=xy-u .
\end{equation}
At $s=u=0$ the ordinary zero is isolated, and its local algebra has basis
\begin{equation}\label{analytic:eq:examplebasis}
 (b_1,\ldots,b_5)=(1,y,y^2,y^3,x),\qquad\mu=5 ,
\end{equation}
since the leading monomials of a local polynomial reduction are
$x^2,xy,y^4$. The defining relations imply
\begin{equation}\label{analytic:eq:examplerelations}
 xy=u,\qquad x^2=s+y^3,\qquad y^4=ux-sy ,
\end{equation}
so with columns representing images of the ordered basis
\eqref{analytic:eq:examplebasis},
\begin{equation}\label{analytic:eq:examplematrices}
 M_x=\begin{pmatrix}
 0&u&0&0&s\\0&0&u&0&0\\0&0&0&u&0\\0&0&0&0&1\\1&0&0&0&0
 \end{pmatrix},\qquad
 M_y=\begin{pmatrix}
 0&0&0&0&u\\1&0&0&-s&0\\0&1&0&0&0\\0&0&1&0&0\\0&0&0&u&0
 \end{pmatrix},
\end{equation}
satisfying exactly
\begin{equation}\label{analytic:eq:matrixidentities}
 M_xM_y=M_yM_x=uI,\qquad M_x^2-M_y^3=sI .
\end{equation}
Starting from the vector for $1$ these matrices generate the five basis
vectors, so they give a faithful five-dimensional realization of the algebra
described by \eqref{analytic:eq:examplerelations};
\cref{analytic:thm:conservationR} supplies the same conclusion for the
analytic monad quotient. Their characteristic polynomials are
\begin{equation}\label{analytic:eq:characteristicexample}
 \chi_y(Y)=Y^5+sY^2-u^2,\qquad \chi_x(X)=X^5-sX^3-u^3 .
\end{equation}
When $u\ne0$ every root of $\chi_y$ is nonzero and determines one solution
by $x=u/y$. In all cases the total local length is five: the polynomial
equations are elimination tools, and a chosen coordinate need not
distinguish all points.

\paragraph{The actual collision equation.}
The Jacobian is
\begin{equation}\label{analytic:eq:exampleJ}
 J_F=\det\begin{pmatrix}2x&-3y^2\\ y&x\end{pmatrix}
 =2x^2+3y^3=2s+5y^3\quad\text{in }A_F ,
\end{equation}
and a direct calculation gives
\begin{equation}\label{analytic:eq:collision}
 \det(M_{J_F})=3125u^6-108s^5 .
\end{equation}
By \cref{analytic:prop:jacobian} the system therefore has five simple
points exactly when the right-hand side is nonzero. This is the collision
equation of the \emph{whole algebra}, not the discriminant of a projected
coordinate: indeed
\begin{equation}\label{analytic:eq:projecteddiscriminant}
 \disc_Y(Y^5+sY^2-u^2)=u^2\,(3125u^6-108s^5),
\end{equation}
and the extra factor $u^2$ is a projection artifact. At $u=0$, $s\ne0$ the
system has the two points $(\pm\sqrt s,0)$ and the three points $(0,y)$ with
$y^3=-s$, all five simple, even though the polynomial in $y$ has a double
root at zero. This is why simultaneous multiplication matrices and the full
Jacobian give better finite-map data than a single eliminant.

A nontrivial collision is explicit. Choose $a\in\m_K^\times$, choose
$\eta\in\C$ with $\eta^2=-3/2$, and put
\[
 s=-\tfrac52a^6,\qquad u=\eta a^5 .
\]
Then $(x,y)=(\eta a^3,a^2)$ is a zero with $J_F=0$, and
\eqref{analytic:eq:collision} vanishes. In the $y$-eliminant this is a
double root, its second derivative there being $15a^6\ne0$: writing
$y=a^2Y$, the eliminant is proportional to $(Y-1)^2(2Y^3+4Y^2+6Y+3)$, whose
greatest common divisor with its derivative is $Y-1$. This specialization
therefore has one point of length two and three other simple points, while
at $s=u=0$ all five units of length are concentrated at the origin.

\paragraph{Three infinitesimal scale regimes.}
Take positive $\alpha,\beta\in\Gamma$ and set $s=t^\alpha$, $u=t^\beta$;
the relative sizes are governed by $5\alpha$ and $6\beta$. All roots are
simple, since in the equal case the two leading coefficients in
\eqref{analytic:eq:collision} subtract to $3125-108=3017\ne0$.

\begin{center}
\begin{tabular}{@{}L{0.22\textwidth}L{0.12\textwidth}L{0.24\textwidth}L{0.24\textwidth}@{}}
\toprule
Regime & Number & $v(y)$ & $v(x)$\\
\midrule
$5\alpha<6\beta$ & $2$ & $\beta-\alpha/2$ & $\alpha/2$\\[0.3em]
                 & $3$ & $\alpha/3$ & $\beta-\alpha/3$\\[0.3em]
$5\alpha>6\beta$ & $5$ & $2\beta/5$ & $3\beta/5$\\[0.3em]
$5\alpha=6\beta$ & $5$ & $2\beta/5$ & $3\beta/5$\\
\bottomrule
\end{tabular}
\end{center}

These assertions are proved directly by rescaling rather than by displaying
a Newton polygon. Suppose $5\alpha<6\beta$. Substituting
$y=t^{\beta-\alpha/2}Y$ in $\chi_y=0$ and dividing by $t^{2\beta}$ gives
\[
 Y^2-1+t^{3\beta-5\alpha/2}Y^5=0,
\]
whose error exponent is positive; each of the two simple ordinary roots
$Y=\pm1$ lifts uniquely to a root in its monad, by degree-one preparation or
by the simple-root case of \cref{analytic:thm:zeros}. Next put
$y=t^{\alpha/3}Y$ and divide by $t^{5\alpha/3}$, obtaining
\[
 Y^5+Y^2-t^{2\beta-5\alpha/3}=0,
\]
whose three nonzero simple leading roots satisfy $Y^3=-1$ and lift to the
second group. The valuations in the two groups differ and together they
account for all five roots; since $x=t^\beta/y$, the stated $x$-valuations
follow. For $5\alpha>6\beta$, put $y=t^{2\beta/5}Y$ and divide by
$t^{2\beta}$ to get $Y^5-1+t^{\alpha-6\beta/5}Y^2=0$, whose five ordinary
fifth roots of unity are simple and lift uniquely. In the equal case the
same substitution gives the ordinary polynomial $Y^5+Y^2-1$, which has no
zero root and is squarefree, its discriminant being $3017$; its five
ordinary roots give the five solutions at that scale. This proves the table
in every divisible ordered value group, not merely for real rational
exponents.

\paragraph{Residues without singular denominators.}
In the basis \eqref{analytic:eq:examplebasis} the total residue functional
is
\begin{equation}\label{analytic:eq:examplelambda}
 \Res_F(1)=\Res_F(y)=\Res_F(y^2)=\Res_F(x)=0,\qquad \Res_F(y^3)=1 .
\end{equation}
One way to identify it first on the simple locus with $u\ne0$ is to use the
$y$-eliminant $q(Y)=Y^5+sY^2-u^2$, for which $q'(y)=yJ_F$ at a root; the
classical coefficient identity for a monic quintic gives
\[
 \sum_{q(y)=0}\frac{y^k}{q'(y)}=0\quad(0\le k<4),
 \qquad \sum_{q(y)=0}\frac{y^4}{q'(y)}=1 ,
\]
and these together with $x=u/y$ give
\eqref{analytic:eq:examplelambda}. The fixed-cycle residue depends
holomorphically on the ordinary parameters near $(0,0)$, so equality on the
open simple locus proves the identities throughout the parameter germ and
hence at every infinitesimal Hahn specialization; they may also be obtained
directly from the ordinary residue expansion
\eqref{analytic:eq:residueexpansion}. The Gram matrix is
\begin{equation}\label{analytic:eq:exampleGram}
 \mathcal G_F=\begin{pmatrix}
 0&0&0&1&0\\0&0&1&0&0\\0&1&0&0&0\\1&0&0&-s&0\\0&0&0&0&1
 \end{pmatrix},\qquad\det\mathcal G_F=1 .
\end{equation}
Unlike \eqref{analytic:eq:collision} this determinant never vanishes: the
residue pairing remains perfect at every collision, including the
length-five point at the origin. By contrast the trace Gram matrix is
$\mathcal G_FM_{J_F}$, so
\[
 \det\bigl(\Tr(M_{b_rb_s})\bigr)_{r,s}=3125u^6-108s^5
\]
does degenerate on the collision locus --- the distinction recorded after
\cref{analytic:cor:residueframe}. Finally
$\Res_F(J_F)=\Res_F(2s+5y^3)=5$, the exact count, and on the simple locus
\begin{equation}\label{analytic:eq:examplecancellation}
 \sum_{F(a)=0}\frac1{J_F(a)}=0,\qquad
 \sum_{F(a)=0}\frac{y(a)^3}{J_F(a)}=1,\qquad
 \sum_{F(a)=0}1=5 .
\end{equation}
For monomial parameters in the scale table each $J_F(a)$ is a nonzero
infinitesimal, so every summand of the first sum has infinite modulus; their
cancellation is nevertheless exact, and the total residue is computed by a
fixed nonsingular matrix rather than by subtracting poorly separated pole
contributions.

\subsection{Beyond polynomial and grid-supported perturbations, over
$\RR_2$}\label{analytic:subsec:nonpolynomial}
The finite example is not the limit of the theorem's scope. Consider the
common-domain entire coefficient families
\begin{align}\label{analytic:eq:nonpolynomial}
 F_1(x,y)&=x^2-y^3-\sum_{m\ge1}t^{\,2-1/m}\exp(m!\,x),\\
 F_2(x,y)&=xy-t^{\sqrt2}-t^3\sin(x+y).\nonumber
\end{align}
Choose a divisible workspace containing the real exponents involved. The
support
\[
 E=\{2-1/m:m\ge1\}\cup\{\sqrt2,3\}
\]
is positive and well ordered; its subset $\{2-1/m\}$ has a finite upper cut
at $2$ while its denominators are unbounded, so the construction is not
restricted to a discrete one-generator grid --- this is the situation
described abstractly in \cref{analytic:rem:initial-segments}. All
coefficient functions are entire, so a common ordinary domain is immediate
despite their very large growth.

\begin{corollary}[A nonpolynomial length-five cluster over $\RR_2$]
\label{analytic:cor:nonpolynomial}
The system \eqref{analytic:eq:nonpolynomial} has finitely many common zeros
in the full surcomplex monad of $(0,0)$, of total local multiplicity five.
Its quotient algebra has the exact basis $(1,y,y^2,y^3,x)$. Every entry of
its two multiplication matrices and of its residue Gram matrix has support
contained in $E^*$, the residue Gram matrix is invertible, and the total
residue of its Jacobian equals five.
\end{corollary}
\begin{proof}
Its ordinary leading system is exactly $(x^2-y^3,xy)$, of multiplicity five.
Apply \cref{analytic:thm:conservationR,analytic:thm:zeros,%
analytic:thm:basechange,analytic:thm:residueduality,analytic:thm:trace}.
Each input to a multiplication or Gram matrix entry is ordinary, giving the
support bound $E^*$.
\end{proof}

This corollary does not assert that the explicit polynomial matrices
\eqref{analytic:eq:examplematrices} remain valid for the nonpolynomial
system; they do not. It asserts that the same \emph{basis size and basis
representatives} survive and that the general normal-form recursion
\eqref{analytic:eq:Wrecursion} computes the new entries. Nor does it assert
that all five roots are simple without computing the new Jacobian
determinant in that algebra. The proof is unaffected by replacing $E$ with a
positive well-ordered set in a higher-rank subgroup of $\No$, or by using
several distinct Archimedean classes of scales: the only summation
requirements are well-ordering and finite contribution at each exponent, and
no ordinary bound on coefficient derivatives is required.

\section{Hypothesis audit, limits, and open questions}
\label{analytic:sec:limits}

\subsection{Hypotheses that cannot be dropped silently}
\label{analytic:subsec:hypotheses}
\paragraph{An isolated complete intersection.}
The leading system must have finite local colength, and exactly $n$
equations in $n$ variables are required. A leading system such as
$(xy,xy)$ in two variables has a positive-dimensional zero set, so no finite
$\mu$ is available; the quotient by the single equation $x$ in two variables
has infinite dimension over the coefficient field, and there is no finite
length to conserve. Arbitrary overdetermined perturbations can \emph{lower}
length: the ideal generated by $(z^2,0)$ has colength two in $\C\{z\}$, and
perturbing it to $(z^2,t^\eta z)$ with $\eta>0$ makes the quotient $K$, of
length one, because $t^\eta$ is a unit in $K$. This contradicts nothing,
since \cref{analytic:thm:finite-deformA} requires exactly $n$ equations with
an isolated complete-intersection leading zero.

\paragraph{Positive perturbation support.}
Replacing $x$ by $x-1$ changes the ordinary leading equation and removes its
zero from the monad of the origin; an error with an exponent below zero can
dominate the leading system entirely. Positive support is the mechanism
making $\id+T$ and $1+\delta h$ invertible with the stated support
certificates. Positivity of the perturbation support fixes the leading
system; changing that system may of course change the invariant.

\paragraph{Well-ordering, not mere positivity.}
See \cref{analytic:rem:positivity}: $\{1/r:r\ge1\}$ is positive but not well
ordered, and the corresponding putative Hahn series is not an element of any
of our rings.

\paragraph{One common ordinary neighbourhood --- where it is needed and
where it is not.}
For every statement proved over $\RR_n$ the fixed domain is not an optional
analytic convenience: it is what makes
\cref{analytic:lem:ordinarykoszul} and
\cref{analytic:thm:residueduality} available at all, and individually
holomorphic coefficient germs with radii tending to zero need not be
coherent data in that sense. Conversely, coefficients may have arbitrarily
large ordinary norms on the common domain; no unmentioned boundedness
hypothesis is used. For every statement proved over $\A_n$ the fixed domain
is \emph{not} available, and its absence is the content of the theorem:
\cref{analytic:ex:no-radius} and
\cref{analytic:subsec:radius-free-example} are elements and systems no
common-domain theorem reaches.

\paragraph{A set-sized workspace before passing to the class.}
Noetherianity and finite-dimensional linear algebra are asserted for the
set-sized rings $\A_n$ and $\RR_n(\Gamma)$. The passage to the full surreal
class is made by the explicit base-change arguments of
\cref{analytic:sec:extension} for one finite cluster at a time, never by
treating a proper class as an ordinary Noetherian ring.

\paragraph{A nonzero value group.}
\Cref{analytic:rem:trivial-group} shows the strong Nullstellensatz fails at
$I=(0)$ when $\Gamma=\{0\}$. The hypothesis is essential rather than
aesthetic, and it enters through
\cref{analytic:prop:faithful} and \cref{analytic:lem:density}.

\paragraph{The chosen analytic categories.}
These results do not classify arbitrary fine-locally analytic functions on
the class plane, and they do not settle the supplied manuscripts' separate
questions about essential singularities, resurgent sectorial summation, or
transcendental all-scale atlases. Those concern compatibility of different
expansion structures rather than a fixed positive-support deformation of an
ordinary finite germ.

\subsection{What is proved, and what is not}
\label{analytic:subsec:not-proved}
The following list is deliberately explicit. Each item is a limitation that
one of the two merged manuscripts stated about its own results, and each is
preserved verbatim in substance.

\paragraph{Status of the whole.}
The assertions are proved from explicitly identified classical inputs. The
proposed contribution is the precise radius-free Hahn-germ package, the
common-domain normal-form and residue package, and the comparison between
them --- not the invention of non-Archimedean Weierstrass theory or residue
duality. No named published conjecture is claimed solved, no exhaustive
novelty certification is asserted, and the proofs have not been formally
verified. Neither a Hahn Weierstrass theorem nor a non-Archimedean
Nullstellensatz is claimed as a new general idea; the coefficient ring and
its support restrictions must be compared, not just theorem titles. We do
not claim that our rings lie outside every broader analytic-structure
framework. In the sources inspected we did not locate the exact combined
statements proved here; we offer them as proposed new results in a precisely
specified setting, and they may also be derivable from a broader existing
framework. This is a proposed original synthesis and formulation, not a
certified first theorem in multivariable non-Archimedean analysis. The
article supplies proofs rather than treating a failure to locate an
identical statement as proof of priority.

\paragraph{Scope discipline.}
No proper-class polynomial ring or proper-class sum is used. Every theorem
about a ring concerns a fixed set-sized $\Gamma$. The support bounds are
support bounds, not real norm estimates, and not a claim that successive
approximations converge in the full surreal fine topology; strong
summability remains the definition of the series, and the estimates
\eqref{analytic:eq:taylor-bound} do not turn every sequence of partial sums
into a convergent sequence.

\paragraph{Per-theorem non-claims.}
\begin{itemize}
\item \Cref{analytic:cor:hypersurface-fiber} is a genuine finite-mapping
statement with infinitesimal parameter values, not only a one-variable
factorization at a single point. It does \emph{not} assert that the roots
can be globally labelled by single-valued analytic functions through a
branch locus.
\item \Cref{analytic:thm:strong-nullA} gives existence of a Nullstellensatz
witness, \emph{not} a uniform bound on the complexity of finding one from
arbitrary support presentations.
\item The splitting \eqref{analytic:eq:splitting} is a choice of linear
splitting, \emph{not} a claim of uniqueness of the quotients $h_j(g)$, and
no ordinary norm bound on it is assumed. The same applies to the
vector-space splittings \eqref{analytic:eq:contraction}: no boundedness or
continuity of them is asserted.
\item The classical facts of
\cref{analytic:subsec:classical-residue} are used on one finite family at a
time. We never assert that all coefficient germs of a Hahn perturbation
admit a common contour or a common analytic representative.
\item \Cref{analytic:eq:hahn-residue} is motivated by expanding each
reciprocal $(f_j+E_j)^{-1}$, but the displayed series, \emph{not} a putative
surcomplex contour integral, is the definition; and in
\cref{analytic:thm:residueduality} the cycle is an ordinary one and is not
treated as a fine-continuous parametrized object in $\SC^n$.
\item The word \emph{finite} in
\cref{analytic:thm:finitemap} has its algebraic meaning throughout. No
claim of properness, compactness, or classical path lifting in the full
surreal fine topology is included.
\item \Cref{analytic:thm:conjugacy} is stated for a complex with homology
only in degree zero. In a complex with homology in several degrees,
positive perturbations can produce a nontrivial differential on the homology
space, and the theorem would have to be modified.
\item The recursion \eqref{analytic:eq:Wrecursion} does \emph{not} assert a
finite-time enumeration algorithm for every set-theoretically possible Hahn
input; the finite-dependence statement at a specified coefficient is
unconditional. More generally, the arguments are algebraic and
support-theoretic rather than a complexity analysis: the set of dependencies
of a fixed coefficient is finite, but an arbitrary ordered group or support
need not have an effective representation from which those dependencies can
be enumerated.
\item \Cref{analytic:thm:zeros} protects total length. It neither requires
nor implies that the number of distinct points stays constant.
\item \Cref{analytic:cor:residueframe} identifies bilinear spaces, not
algebras; it does not erase the nonreduced algebraic structure of a
collision, and a perfect residue pairing is not the usually degenerate trace
pairing of a nonreduced finite algebra. At a multiple point
\eqref{analytic:eq:simple-residue} is not a valid formula with a zero
denominator (\cref{analytic:rem:localization}).
\item \Cref{analytic:cor:nonpolynomial} does not assert that the explicit
matrices \eqref{analytic:eq:examplematrices} remain valid for the
nonpolynomial system, nor that all five of its roots are simple without
computing the new Jacobian determinant.
\item No Noether normalization or finite-projection theorem is claimed for
$\A_n$ (\cref{analytic:rem:no-normalization-A}), and no analogue of
\cref{analytic:thm:formal-comparison} is claimed for $\RR_n$.
\item The relative version of the normal-form construction over a
positive-dimensional ordinary parameter space, as a module over the
coherent parameter algebra rather than as complex-linear splittings on the
total space, is \emph{not} claimed: the proofs above deliberately do not
assert that stronger relative linearity.
\item The attainment of both valuation bounds in the Hensel--Rouch\'e
theorem, including the second-order bound, is proved in the companion
report \cite{FiniteDeformations} and is not re-proved here
(\cref{analytic:rem:two-certificates}).
\item The symbolic checks accompanying the source manuscripts concern the
displayed formulas only. They do not implement arbitrary Hahn fields and
are not machine verification of the Noetherian, flatness, residue,
conservation, or stability proofs.
\end{itemize}

\subsection{Specific research directions left open}
\label{analytic:subsec:open}
Four questions are left open by these proofs, and they are questions, not
claims about the current status of any named conjecture.

First, whether natural common-domain versions of the radius-free
construction support useful global gluing and coherent sheaf theorems
without reinstating a hidden uniform-radius assumption. The Noetherian
results proved here are ring theorems on an entire monad, not gluing
theorems, and neither ring is given a sheaf theory, Cartan's theorems for
surcomplex domains, or an extension of ordinary contour integration to
arbitrary fine-topological paths.

Second, a weighted version of
\cref{analytic:thm:local-stability}: the single scalar loss $\kappa_a$
treats every coordinate with the same valuation threshold, and diagonal
weights adapted to the Newton data may improve separation bounds for
anisotropic clusters such as
\eqref{analytic:eq:four-system}. \Cref{analytic:ex:sharpness} prevents a
universal improvement of the scalar factor two but does not preclude
stronger coordinate-sensitive hypotheses.

Third, a detailed comparison with existing separated analytic structures:
one should determine the exact axiomatic status of the radius-free
convergent-coefficient subsystem, including uniformity properties not used
here. \Cref{analytic:thm:formal-comparison} does not by itself prove every
model-theoretic or global analytic property of a larger full separated
system.

Fourth, replacing an isolated complete intersection by a finite
Cohen--Macaulay module with a selected free resolution. The positive
conjugacy argument of \cref{analytic:thm:conjugacy} already applies to any
complex with homology only in degree zero, provided a perturbation of its
differential is supplied that still squares to zero; the genuinely new
problem is to construct the required compatible perturbation of \emph{all}
syzygies rather than of the first equations only. For complete
intersections the Koszul construction supplies it automatically.

\appendix
\section{The finite-family classical input and the trace proof}
\label{analytic:app:classical-family}

We spell out the ordinary deformation statement invoked in
\cref{analytic:thm:trace}. Let
\[
 g_i(z,u)=f_i(z)+\sum_{\ell=1}^qu_\ell q_{i\ell}(z),
 \qquad q_{i\ell}\in R_n,
\]
with $f$ satisfying \eqref{analytic:eq:mu}. The germs in this \emph{finite}
family have a common ordinary representative neighbourhood --- which is
exactly why the reduction of \cref{analytic:lem:finite-dependency} is what
makes the radius-free case work. The ordinary finite-mapping theorem,
obtained from Weierstrass division, says that
$C=\C\{z,u\}/(g_1,\ldots,g_n)$ is a finite module over $\C\{u\}$ near the
distinguished fibre. The sequence consisting of the $g_i$ followed by the
parameters $u_\ell$ is a system of parameters in the regular local ring
$\C\{z,u\}$, since quotienting by it gives the finite-dimensional algebra
$R_n/(f)$; regular local rings are Cohen--Macaulay, so this system of
parameters is a regular sequence, and hence the parameters form a regular
sequence on $C$. The usual finite-module flatness criterion over the regular
parameter ring shows $C$ is free over $\C\{u\}$ of rank $\mu$, and its
reduction modulo $(u)$ has basis the chosen $b_1,\ldots,b_\mu$, so those
same lifts form a basis by Nakayama's lemma. These are standard local
analytic and commutative-algebra facts
\cite{GrauertRemmert,GunningRossi,Stacks}. Equivalently, the relative zero
space is finite over the parameter domain, is a complete intersection and
hence Cohen--Macaulay, and is therefore flat over the smooth parameter base.

For small representatives the finite algebra records the entire cluster near
zero, not merely a separately chosen point of a nearby fibre. Classical
local duality equips it with a parameter-holomorphic residue functional, and
expanding the reciprocal factors on an ordinary protecting residue cycle
gives
\[
 \res^{\mathrm{cluster}}_{g_u}(H)
 =\sum_{k\in\N^n}(-1)^{|k|}\res_{f^{k+1}}
 \Bigl(H\prod_i\Bigl(\sum_\ell u_\ell q_{i\ell}\Bigr)^{k_i}\Bigr),
\]
an ordinary convergent expansion for sufficiently small parameters;
equivalently it is an equality of formal parameter series, which is all the
Hahn proof needs. Powers of the $f_i$ in the denominators do not change the
residue cycle class used to extract the coefficients.

The classical trace--Jacobian identity may be understood first on fibres
where all points are simple, using a change of variables at each point and
the decomposition of the algebra into one-dimensional factors. To include
multiple fibres, augment the family by independent constant parameters in
the equations; generic small target values of a finite holomorphic map in
characteristic zero are regular values, and both the cluster residue and the
trace matrix are holomorphic in all parameters, so equality on the regular
fibres extends to all fibres and one then sets the added parameters to
zero. The algebraic version is the theorem of \cite{SS}; related finite
intersection deformations appear in \cite{KW}.

There is a subtle point about the arbitrary splitting $h$ in
\eqref{analytic:eq:splitting}, and about the contraction
\eqref{analytic:eq:contraction}. The Neumann series built from either need
not converge in an ordinary analytic norm. Nevertheless, in the $(u)$-adic
formal completion it gives a valid identity: every coefficient of total
parameter degree below $N$ uses finitely many Neumann terms, and since the
completed finite algebra is free with the chosen basis those formal
coordinates are unique and coincide with the Taylor coefficients of the
actual analytic coordinates. This is why no continuity assertion about $h$
is hidden in \cref{analytic:thm:trace}. Uniqueness also follows directly
from \cref{analytic:thm:conjugacy}, using ordinary parameter degree to form
the formal inverse.

Finally, the Hahn application never needs an infinite-parameter version of
this ordinary argument. At a chosen exponent,
\cref{analytic:lem:finite-dependency} reduces everything to one of these
finite families and a finite list of its coefficient identities. That
reduction is the precise substitute for the common contour that the original
infinite collection of germs need not possess. No interchange of a
fine-topological limit with an integral is involved anywhere.

\section{Support accounting and proof dependencies}
\label{analytic:app:support}

The dependency chain, in both settings, is
\begin{gather*}
 \text{Hahn--Neumann lemma}\ \Longrightarrow\
 \text{positive operator inversion},\\
 \text{positive operator inversion}\ \Longrightarrow\
 \text{division and preparation},\\
 \text{positive operator inversion}\ \Longrightarrow\
 \text{Koszul conjugacy}\ \Longrightarrow\
 \text{fixed-basis conservation}.
\end{gather*}
Division supplies the Noetherian and point-identification route in each
ring. Classical residue theory supplies the leading perfect form and the
finite-parameter trace identity, and the support calculus gives the Hahn
extension. The order avoids circularity: division first proves
Noetherianity, then the maximal-ideal theorem and the strong
Nullstellensatz; the complete-intersection normal form proves spanning
without knowing the perturbed dimension; the residue is defined
independently of roots and annihilates the perturbed ideal; its ordinary
leading Gram matrix proves independence and perfect duality; only then are
maximal ideals interpreted as points and local dimensions summed. The trace
theorem uses ordinary finite-family residue theory, not the
still-to-be-proved Hahn zero count, and the local stability theorem uses the
already established length-one deformation case after a rigorously allowed
rescaling.

In the table, $S$ is a source support, $E$ a well-ordered positive
perturbation support, and $E_a$ the union of the supports of an
infinitesimal displacement.

\begin{center}\small
\begin{longtable}{@{}L{0.40\textwidth}L{0.52\textwidth}@{}}
\toprule
Construction & Containing exponent set, or controlling finite-dimensional
fact\\
\midrule
\endhead
Coefficientwise linear operation or differentiation & A subset of the input
support $S$\\[0.35em]
Evaluation at $c+\varepsilon$, or infinitesimal translation &
$S+E_a^*$\\[0.35em]
Positive-support operator inverse on an input & $S+E^*$\\[0.35em]
Division quotient and remainder, over $\A_n$ or over $\RR_n$ &
$S+E^*$\\[0.35em]
Normalized preparation factors & $E^*$\\[0.35em]
Koszul inverse $(1+\delta h)^{-1}$, over $\RR_n$ & $S+E^*$, on one ordinary
domain\\[0.35em]
Complete-intersection normal coordinates and division certificates &
$S+E^*$\\[0.35em]
Perturbed residue of a numerator & $S+E^*$\\[0.35em]
Residue Gram matrix, its inverse, and the constant frame $S$ & $E^*$; the
leading Gram determinant is nonzero\\[0.35em]
Multiplication constants of the fixed ordinary basis & $E^*$; reduction
gives the ordinary nilpotent matrices\\[0.35em]
Individual root coordinates & Algebraic over $K$; \emph{not} necessarily
supported in $E^*$\\[0.35em]
Workspace enlargement & Same finite basis, same matrices, algebraic
closedness of the original $K$\\
\bottomrule
\end{longtable}
\end{center}

The support lemma proves two things for every entry: the containing set is
well ordered, and only finitely many source terms contribute at one
exponent. The second is what licenses regrouping, coefficientwise
differentiation, and comparison of a chosen coefficient with a
finite-parameter problem. A mere lower bound on valuation would provide no
such licence; this is the content of
\cref{analytic:rem:rank,analytic:rem:initial-segments}.

\section{Notation}\label{analytic:app:notation}

\begin{center}\small
\begin{longtable}{@{}L{0.24\textwidth}L{0.68\textwidth}@{}}
\toprule
Symbol & Meaning\\
\midrule
\endhead
$\SC=\No[i]$ & Surcomplex class field\\[0.3em]
$\Gamma$, $t^\gamma=\omega^{-\gamma}$ & Nonzero divisible set-sized value
group; normal-form monomial, positive exponent meaning infinitesimal\\[0.3em]
$K=K_\Gamma$, $K^\circ$, $\m_K$ & $\C((t^\Gamma))$, its valuation ring, and
its ideal of infinitesimals\\[0.3em]
$v$, $\vH$, $\vmin$, $\red$ & Field valuation ($v(0)=+\infty$); Hahn
valuation on a coefficient ring; minimum coordinate valuation;
reduction\\[0.3em]
$R_n=\C\{z\}$ & Ordinary convergent germs at the origin\\[0.3em]
$\OO(D)((t^\Gamma))$ & Fixed-polydisk ring\\[0.3em]
$\HH_\Gamma(U)$, $\RR_n$ & Hahn families holomorphic on one $U$; their germ
ring at the ordinary origin (common-domain ring)\\[0.3em]
$\A_n=R_n((t^\Gamma))$ & Radius-free Hahn-germ ring\\[0.3em]
$\FF_n=\C[[z]]((t^\Gamma))$ & Formal-coefficient Hahn ring\\[0.3em]
$E$, $E^*$ & Positive well-ordered perturbation support; the additive monoid
it generates, including the empty sum\\[0.3em]
$\m_a$ & Evaluation ideal $(z_1-a_1,\ldots,z_n-a_n)$\\[0.3em]
$\mu$, $e_a$ & Global colength; local algebraic multiplicity, equal to the
formal multiplicity at the actual displaced zero\\[0.3em]
$B_0$, $B_F$, $A_F$ & Ordinary local algebra; deformed algebra over $\A_n$;
deformed algebra over $\RR_n$\\[0.3em]
$d_f$, $\delta$, $h$, $U$ & Ordinary Koszul differential, its positive
perturbation, a fixed contraction, and $1+\delta h$\\[0.3em]
$\NF_F$, $Q_j$, $M_H$ & Basis-coordinate normal form; division certificates;
multiplication matrix\\[0.3em]
$\res_f$, $\Res_F$, $\Lambda_F$, $\mathcal G_F$ & Ordinary local residue;
the Hahn series residue; the contour-expansion residue over $\RR_n$, equal
to $\Res_F$ where both are defined; the residue Gram matrix\\[0.3em]
$J_F$, $\kappa_a$, $\mathcal B_a$ & Jacobian determinant; inverse-Jacobian
valuation loss; stability ball\\
\bottomrule
\end{longtable}
\end{center}

\begin{thebibliography}{99}
\addcontentsline{toc}{section}{References}

\bibitem{Alling}
Norman L. Alling,
\emph{Foundations of Analysis over Surreal Number Fields},
North-Holland Mathematics Studies \textbf{141}, North-Holland, 1987.

\bibitem{BM}
Alessandro Berarducci and Vincenzo Mantova,
``Transseries as germs of surreal functions,''
\emph{Transactions of the American Mathematical Society} \textbf{371}
(2019), 3549--3592.
\href{https://arxiv.org/abs/1703.01995}{arXiv:1703.01995}.

\bibitem{CDS}
Eduardo Cattani, Alicia Dickenstein, and Bernd Sturmfels,
``Computing multidimensional residues,'' in L.~Gonz\'alez-Vega and
T.~Recio (eds.), \emph{Algorithms in Algebraic Geometry and Applications},
Progress in Mathematics \textbf{143}, Birkh\"auser, 1996, pp.~135--164.
\href{https://arxiv.org/abs/alg-geom/9404011}{arXiv:alg-geom/9404011};
\href{https://doi.org/10.1007/978-3-0348-9104-2_8}{doi:10.1007/978-3-0348-9104-2\_8}.
Sections 0 and 4 give the residue-duality and trace background used here.

\bibitem{CL}
Raf Cluckers and Leonard Lipshitz,
``Fields with analytic structure,''
\emph{Journal of the European Mathematical Society} \textbf{13} (2011),
1147--1223.
\href{https://doi.org/10.4171/JEMS/278}{doi:10.4171/JEMS/278};
\href{https://arxiv.org/abs/0908.2376}{arXiv:0908.2376}.
The arbitrary-support, common-radius coefficient families are in \S3.1,
Remark 3.1.4 and Example 3.3(6) of the arXiv text; related normalization and
Nullstellensatz results are in \S5.2, and Propositions 5.2.5--5.2.7 are
close antecedents of the structure theorems proved here.

\bibitem{CLR}
Raf Cluckers, Leonard Lipshitz, and Zachary Robinson,
``Real closed fields with nonstandard and standard analytic structure,''
\emph{Journal of the London Mathematical Society} (2) \textbf{78} (2008),
198--212.
\href{https://doi.org/10.1112/jlms/jdn024}{doi:10.1112/jlms/jdn024};
\href{https://arxiv.org/abs/math/0610666}{arXiv:math/0610666}.
Theorem 2.5 and Remark 2.6 give multivariable preparation and division for
nonstandard analytic structures.

\bibitem{CompanionA}
\emph{Surcomplex Analysis: Infinitesimal Calculus, Hahn-Coherent Holomorphy,
and Contour Theory},
supplied research manuscript, September 21, 2026.
Source filename: \texttt{surcomplex\_analysis(3).tex}.
Its several-variable research target motivates this continuation; none of
its unproved targets is used as a theorem.

\bibitem{CompanionB}
\emph{Surcomplex Analysis: Infinitesimal Calculus, Hahn-Coherent Holomorphy,
and Contour Theory},
supplied unpublished manuscript, September 21, 2026.
Source filename: \texttt{surcomplex\_analysis(2).tex}.
The present report uses its Hahn-coherent framework and addresses its
concluding several-variable research direction. No authorship or published
priority is inferred from the file.

\bibitem{CE}
Ovidiu Costin and Philip Ehrlich,
``Integration on the surreals,''
\href{https://arxiv.org/abs/2208.14331}{arXiv:2208.14331}, version 5, 2024.

\bibitem{Crainic}
Marius Crainic,
``On the perturbation lemma, and deformations,'' 2004.
\href{https://arxiv.org/abs/math/0403266}{arXiv:math/0403266}.
Our contraction convention is $dh+hd=1-ip$, and
\cref{analytic:thm:conjugacy} is proved with that sign convention.

\bibitem{FiniteDeformations}
\emph{Finite Hahn Deformations: Division, Conservation of Multiplicity,
and Residue Duality}, companion merged report in this collection,
\texttt{docs/surcomplex/finite-deformations/}, September 21, 2026.
It states the deformation package at its strongest hypotheses --- arbitrary
isolated complete intersection, no simple zero, no triangular system, no
polynomial perturbation, no common radius --- together with both stability
certificates and the attainment of the valuation bounds.
Its Section 15.3 is the statement of record for the lossiness of the
determinant-only threshold, $0\leq\kappa_a\leq v(\det A)$, which the present
report cites rather than restates.

\bibitem{Contours}
\emph{Jordan Separation, Cauchy Theory, and Stokes Calculus on the
Surcomplex Plane}, companion merged report in this collection,
\texttt{docs/surcomplex/contours-and-stokes/}, September 21, 2026.
It addresses the separate question of supplying a surcomplex contour that
represents the series \eqref{analytic:eq:hahn-residue}.

\bibitem{GrauertRemmert}
Hans Grauert and Reinhold Remmert,
\emph{Coherent Analytic Sheaves},
Grundlehren der mathematischen Wissenschaften \textbf{265}, Springer, 1984.
Classical Weierstrass division, local analytic algebras, and finite
mappings.

\bibitem{GH}
Phillip Griffiths and Joseph Harris,
\emph{Principles of Algebraic Geometry}, Wiley, 1978; Wiley Classics Library
edition, 1994, Chapter 5, ``Residues.''
\href{https://doi.org/10.1002/9781118032527}{doi:10.1002/9781118032527}.

\bibitem{GunningRossi}
Robert C. Gunning and Hugo Rossi,
\emph{Analytic Functions of Several Complex Variables},
Prentice-Hall, 1965; reprinted by AMS Chelsea Publishing, 2009.
Ordinary Weierstrass theory, coherent analytic sheaves, and Stein-space
theorems, including Cartan's theorem~B as used in
\cref{analytic:lem:ordinarykoszul}.

\bibitem{KW}
Ernst Kunz and Rolf Waldi,
``Deformations of zero-dimensional intersection schemes and residues,''
\emph{Note di Matematica} \textbf{11} (1991), 247--259.

\bibitem{Neumann}
B. H. Neumann,
``On ordered division rings,''
\emph{Transactions of the American Mathematical Society} \textbf{66} (1949),
202--252.
Support and well-ordering results for generalized series.

\bibitem{Poonen}
Bjorn Poonen,
``Maximally complete fields,''
\emph{L'Enseignement Math\'ematique} (2) \textbf{39} (1993), 87--106.
\href{https://math.mit.edu/~poonen/papers/amsval.pdf}{Author's full text}.
Corollary 4 gives the algebraic closure result for divisible value group and
algebraically closed residue field.

\bibitem{SS}
G\"unter Scheja and Uwe Storch,
``\"Uber Spurfunktionen bei vollst\"andigen Durchschnitten,''
\emph{Journal f\"ur die reine und angewandte Mathematik} \textbf{278/279}
(1975), 174--190.
\href{https://doi.org/10.1515/crll.1975.278-279.174}{doi:10.1515/crll.1975.278-279.174}.

\bibitem{Stacks}
The Stacks Project Authors, \emph{The Stacks Project}.
\href{https://stacks.math.columbia.edu/tag/00FZ}{Section 10.35, Jacobson
rings (tag 00FZ)};
\href{https://stacks.math.columbia.edu/tag/00MC}{Lemma 10.97.3, faithful
flatness of completion (tag 00MC)};
\href{https://stacks.math.columbia.edu/tag/00GQ}{tag 00GQ} for lying over;
and the chapter ``Commutative Algebra'' for finite and integral extensions,
regular local rings, Cohen--Macaulay rings and flatness. Consulted
September 21, 2026.

\bibitem{Tsikh}
A. K. Tsikh,
\emph{Multidimensional Residues and Their Applications},
Translations of Mathematical Monographs \textbf{103}, American Mathematical
Society, 1992.
\href{https://doi.org/10.1090/mmono/103}{doi:10.1090/mmono/103}.

\end{thebibliography}
\end{document}
